\section{Proofs from Section~\ref{subsec:mean-field}}
\label{sec:proof-mean-field}

Throughout this section, the hypotheses of
Theorem~\ref{theorem:trajectory} hold.
We first establish well-definedness of the recursion in
Definition~\ref{def:mean-field-recursion}
in Section~\ref{subsec:well-definedness-effective}.
We introduce notation and state the master recursion in
Section~\ref{subsec:notation-master-rec}.
After proving the supporting lemmas in
Section~\ref{subsec:supporting-lem}, we establish the master
recursion in Section~\ref{subsec:proof-master-rec}
and deduce the trajectory theorem and its corollary in
Section~\ref{subsec:proof-trajectory}.
Fix $T\in\mathbb N$.

\subsection{Well-definedness of the effective process}
\label{subsec:well-definedness-effective}

For $a\in\{w,g\}$, write
\[
    \bar C_t^{[a]}
    :=
    \bigl(\bar C^{[a]}(s,u)\bigr)_{0\leq s,u\leq t},
    \qquad
    \bar\alpha_t^{[a]}
    :=
    \bigl(\bar\alpha^{[a]}(t,s)\bigr)_{0\leq s<t}.
\]
We also write
\[
    \bar\phi_t^{[w]}
    :=
    \bigl(\bar\phi^{[w]}(t,s)\bigr)_{0\leq s<t},
    \qquad
    \bar\phi_t^{[g]}
    :=
    \bigl(\bar\phi^{[g]}(t,s)\bigr)_{0\leq s\leq t}.
\]
Coefficient vectors are column vectors, and matrices and
vectors with empty index sets are understood to be empty.

\begin{lemma}[Well-definedness and positive innovations]
        \label{lemma:mf-proof-nondegeneracy}
        Under the hypotheses of Theorem~\ref{theorem:trajectory},
        the effective recursion in
        Definition~\ref{def:mean-field-recursion}
        uniquely determines the sample-side
        and parameter-side joint trajectory laws at every finite
        horizon, and all its coordinates have finite second moments. For every $t\geq0$,
        \[
            \bar C_t^{[w]}\succ0,
            \qquad
            \bar C_t^{[g]}\succ0,
            \qquad
            \bar v_t^{[w]}>0,
            \qquad
            \bar v_t^{[g]}>0.
        \]
        Moreover, $\tilde r^t$ has an absolutely continuous distribution and
        \[
            \tilde{\mathbb P}
            (\tilde r^t\in\{0,\kappa_-,\kappa_+\})=0.
        \]
        \end{lemma}

        \begin{proof}
        Fix the base variables and independent Gaussian innovation
        sequences specified in
        Definition~\ref{def:mean-field-recursion}.
        We construct the processes in the order
        \[
            \tilde w^t
            \ \longrightarrow\
            (\tilde q^t,\tilde r^t,\tilde g^t)
            \ \longrightarrow\
            (\tilde p^t,\tilde w^{t+1}).
        \]
        We use the following one-step forms of the effective fluctuations:
        \begin{align}
            \tilde q^t
            &=
            \sum_{s=0}^{t-1}
                \bar\alpha^{[w]}(t,s)\tilde q^s
            +
            \frac1{\sqrt{\bar\delta}}
            \sum_{s=0}^{t-1}
                \bar\phi^{[w]}(t,s)\tilde g^s
            +
            \bar v_t^{[w]}z_t^{[g]},
            \label{eq:mf-proof-forward-one-step}\\
            \tilde p^t
            &=
            \sum_{s=0}^{t-1}
                \bar\alpha^{[g]}(t,s)\tilde p^s
            +
            \sqrt{\bar\delta}
            \sum_{s=0}^{t}
                \bar\phi^{[g]}(t,s)\tilde w^s
            +
            \bar v_t^{[g]}z_t^{[w]}.
            \label{eq:mf-proof-backward-one-step}
        \end{align}
        Whenever the coefficients involved have been constructed,
        these identities follow from the cumulative-memory definitions:
        substituting the representations of the previous fluctuations
        and collecting coefficients gives precisely the defining
        recursions for $\bar\gamma^{[w]}$, $\bar\gamma^{[g]}$,
        $\xi^{[g]}$, and $\xi^{[w]}$.
        
        \begin{enumerate}[wide]
            \item \emph{Initialization.}
            The initialization assumption gives
            \[
                \tilde{\mathbb E}
                [\tilde\mu^2+(\tilde w^0)^2]<\infty,
                \qquad
                \bar b^0\in\mathbb R,
            \]
            and
            \[
                \bar C_0^{[w]}
                =
                \tilde{\mathbb E}[(\tilde w^0)^2]
                =
                (\bar v_0^{[w]})^2
                >0.
            \]
            At $t=0$, coefficient systems and sums involving
            earlier times are empty.
        
            \item \emph{Forward construction and residual nondegeneracy.}
            Suppose that the parameter variables have been constructed
            through $\tilde w^t$ and $\tilde p^{t-1}$, and the sample
            variables through time $t-1$, all with finite second moments.
            Suppose also that $\bar b^t$ and the previously constructed
            deterministic coefficients are finite, and that
            \[
                \bar C_t^{[w]}\succ0,
                \qquad
                \bar C_{t-1}^{[g]}\succ0
                \quad\text{when }t\geq1.
            \]
            The systems defining $\bar\alpha_t^{[w]}$ and
            $\bar\phi_t^{[w]}$ have unique solutions:
            their coefficient matrices are
            $\bar C_{t-1}^{[w]}$ and $\bar C_{t-1}^{[g]}$,
            respectively.
            All entries on their right-hand sides are finite
            by Cauchy--Schwarz.
            The projection normal equations give
            \[
                (\bar v_t^{[w]})^2
                =
                \inf_{a\in\mathbb R^t}
                \tilde{\mathbb E}
                \left[
                    \left(
                        \tilde w^t-\sum_{s=0}^{t-1}a_s\tilde w^s
                    \right)^2
                \right].
            \]
            For $t\geq1$, this quantity is strictly positive:
            otherwise $\tilde w^t$ would belong to the linear span
            of its predecessors in $L^2$, contradicting
            $\bar C_t^{[w]}\succ0$.
            Positivity at $t=0$ follows from the initialization.
            The forward recursion therefore defines $\tilde q^t\in L^2$.
            Also,
            $\bar m^t=\tilde{\mathbb E}[\tilde\mu\tilde w^t]$
            is finite, so the score equation defines $\tilde r^t\in L^2$.
        
            On the sample-side space, set
            \[
                \mathcal H_t^{[g]}
                :=
                \sigma(
                    \tilde Y,\tilde\Delta,
                    z_0^{[g]},\ldots,z_{t-1}^{[g]}
                ).
            \]
            The one-step forward identity and the score equation give
            \begin{equation}
                \tilde r^t
                =
                A_t+\sigma\bar v_t^{[w]}z_t^{[g]},
                \qquad
                A_t\in L^2(\mathcal H_t^{[g]}),
                \label{eq:mf-proof-score-fresh-noise}
            \end{equation}
            where $z_t^{[g]}$ is independent of $\mathcal H_t^{[g]}$.
            Since $\sigma\bar v_t^{[w]}>0$,
            \[
                \mathscr L(\tilde r^t\mid\mathcal H_t^{[g]})
                =
                \mathcal N\bigl(A_t,\sigma^2(\bar v_t^{[w]})^2\bigr).
            \]
            In particular, for every Lebesgue-null Borel set
            $B\subseteq\mathbb R$,
            \[
                \tilde{\mathbb P}(\tilde r^t\in B)
                =
                \tilde{\mathbb E}
                \left[
                    \tilde{\mathbb P}
                    (\tilde r^t\in B\mid\mathcal H_t^{[g]})
                \right]
                =0.
            \]
            Thus $\tilde r^t$ has an absolutely continuous distribution,
            and it assigns zero probability to
            $\{0,\kappa_-,\kappa_+\}$.
            The Lipschitz assumptions imply
            \[
                |\ell'(y,r)|
                \leq
                \max_{y'\in\{-1,+1\}}|\ell'(y',0)|
                +L_\ell|r|.
            \]
            Since the selection indicators are bounded and
            $\bar\rho\in(0,1)$, there is a deterministic constant
            $C<\infty$ such that
            \[
                |\tilde g^t|\leq C(1+|\tilde r^t|).
            \]
            Consequently, $\tilde g^t\in L^2$.
            We next prove positivity of its innovation.
            Since $\tilde{\mathbb P}(\tilde\Delta=1)=\bar\rho>0$,
            there exists $y^\star\in\{-1,+1\}$ such that
            \[
                E^\star
                :=
                \{\tilde Y=y^\star,\tilde\Delta=1\}
                \quad\text{satisfies}\quad
                \tilde{\mathbb P}(E^\star)>0.
            \]
            This event belongs to $\mathcal H_t^{[g]}$, and on it
            \[
                \tilde g^t
                =
                \frac1{\bar\rho}\ell'(y^\star,\tilde r^t).
            \]
            Conditionally on almost every history in $E^\star$,
            the score has a Gaussian density that is strictly positive
            on $\mathbb R$.
            Since $\ell'(y^\star,\cdot)$ is continuous and nonconstant,
            its value under this Gaussian law cannot be almost surely
            constant. Indeed, such constancy would imply that the
            function is constant Lebesgue-almost everywhere, and
            continuity would then make it constant everywhere.
            Hence
            \[
                \operatorname{Var}
                (\tilde g^t\mid\mathcal H_t^{[g]})>0
                \quad\text{almost surely on }E^\star,
            \]
            and therefore
            \[
                \tilde{\mathbb E}
                \left[
                    \operatorname{Var}
                    (\tilde g^t\mid\mathcal H_t^{[g]})
                \right]>0.
            \]
            The previous residuals are $\mathcal H_t^{[g]}$-measurable.
            For any $a\in\mathbb R^t$, the conditional square expansion
            gives
            \[
            \begin{aligned}
                \tilde{\mathbb E}
                \left[
                    \left(
                        \tilde g^t-\sum_{s=0}^{t-1}a_s\tilde g^s
                    \right)^2
                \right]
                =
                \tilde{\mathbb E}
                \left[
                    \operatorname{Var}
                    (\tilde g^t\mid\mathcal H_t^{[g]})
                \right]
                +
                \tilde{\mathbb E}
                \left[
                    \left(
                        \tilde{\mathbb E}
                        [\tilde g^t\mid\mathcal H_t^{[g]}]
                        -\sum_{s=0}^{t-1}a_s\tilde g^s
                    \right)^2
                \right].
            \end{aligned}
            \]
            The residual projection is well defined because
            $\bar C_{t-1}^{[g]}$ is positive definite when $t\geq1$.
            Taking the infimum over $a$ yields
            \begin{equation}
                (\bar v_t^{[g]})^2
                \geq
                \tilde{\mathbb E}
                \left[
                    \operatorname{Var}
                    (\tilde g^t\mid\mathcal H_t^{[g]})
                \right]
                >0.
                \label{eq:mf-proof-residual-positive-innovation}
            \end{equation}
            It follows that $\bar C_t^{[g]}\succ0$.
            Indeed, a vanishing linear combination of
            $\tilde g^0,\ldots,\tilde g^t$ with a nonzero coefficient
            of $\tilde g^t$ would contradict
            \eqref{eq:mf-proof-residual-positive-innovation}.
            A nontrivial combination with zero last coefficient
            is excluded by $\bar C_{t-1}^{[g]}\succ0$.
            At $t=0$, positivity follows directly from
            $(\bar v_0^{[g]})^2>0$.
        
            Finally,
            $\bar\zeta^t=\tilde{\mathbb E}[\tilde g^t]$ and
            $\bar\chi^t=\tilde{\mathbb E}[\tilde Y\tilde g^t]$
            are finite, and so is
            $\bar b^{t+1}=\bar b^t-\eta\bar\zeta^t$.
        
            \item \emph{Backward construction and the next weight innovation.}
            The residual projection coefficients have already been
            determined, and the system defining $\bar\phi_t^{[g]}$
            has coefficient matrix $\bar C_t^{[w]}\succ0$.
            Its right-hand side is finite by Cauchy--Schwarz.
            Hence the backward coefficients are uniquely determined,
            and the backward recursion defines $\tilde p^t\in L^2$.
        
            On the parameter-side space, set
            \[
                \mathcal H_t^{[w]}
                :=
                \sigma(
                    \tilde\mu,\tilde w^0,
                    z_0^{[w]},\ldots,z_{t-1}^{[w]}
                ).
            \]
            All weights through time $t$ and backward fluctuations
            through time $t-1$ are measurable with respect to this
            history.
            The one-step backward identity and the parameter update
            therefore give
            \begin{equation}
                \tilde w^{t+1}
                =
                B_t
                -
                \frac{\eta\sigma}{\sqrt{\bar\delta}}
                \bar v_t^{[g]}z_t^{[w]},
                \qquad
                B_t\in L^2(\mathcal H_t^{[w]}).
                \label{eq:mf-proof-weight-fresh-noise}
            \end{equation}
            Here square integrability follows from the preceding
            moment bounds and
            \[
                |j'(u)|\leq |j'(0)|+L_j|u|.
            \]
            Since $z_t^{[w]}$ is independent of $\mathcal H_t^{[w]}$,
            \[
                \operatorname{Var}
                (\tilde w^{t+1}\mid\mathcal H_t^{[w]})
                =
                \frac{\eta^2\sigma^2}{\bar\delta}
                (\bar v_t^{[g]})^2.
            \]
            The next weight projection is well defined because its
            coefficient matrix is $\bar C_t^{[w]}\succ0$.
            Applying the same conditional square expansion as above,
            now to linear combinations of
            $\tilde w^0,\ldots,\tilde w^t$, yields
            \begin{equation}
            \begin{aligned}
                (\bar v_{t+1}^{[w]})^2
                =
                \inf_{a\in\mathbb R^{t+1}}
                \tilde{\mathbb E}
                \left[
                    \left(
                        \tilde w^{t+1}
                        -\sum_{s=0}^{t}a_s\tilde w^s
                    \right)^2
                \right]
                \geq
                \tilde{\mathbb E}
                \left[
                    \operatorname{Var}
                    (\tilde w^{t+1}\mid\mathcal H_t^{[w]})
                \right]
                =
                \frac{\eta^2\sigma^2}{\bar\delta}
                (\bar v_t^{[g]})^2
                >0.
            \end{aligned}
            \label{eq:mf-proof-weight-positive-innovation}
            \end{equation}
            The same linear-independence argument used for the
            residuals proves $\bar C_{t+1}^{[w]}\succ0$.
            Thus the induction hypotheses hold at time $t+1$.
        \end{enumerate}
        
        Starting from the initialization, this alternating construction
        establishes the stated moment bounds, positive definiteness,
        positive innovations, and absolute continuity at every finite time.
        The pseudo-label variables are bounded, and the remaining
        Gaussian and memory variables are finite linear combinations
        of variables already shown to have finite second moments.
        
        At every half-step, the coefficient systems have unique solutions,
        their moments are determined by the previously constructed laws,
        and the new variables are specified functions of the existing
        variables and fresh Gaussian innovations.
        Thus the recursion is uniquely determined once the base variables
        and innovation sequences are fixed, and its sample-side and
        parameter-side trajectory laws are uniquely determined.
        \end{proof}

\subsection{Notation and master recursion}
\label{subsec:notation-master-rec}

We first rewrite Algorithm~\ref{alg:self-training} in matrix form
to isolate its dependence on the Gaussian design noise.
Let $X$ and $Y$ denote the design matrix and label vector
\[
    X
    :=
    \begin{bmatrix}
        X_1^\top\\
        \vdots\\
        X_n^\top
    \end{bmatrix}
    \in\mathbb R^{n\times d},
    \qquad
    Y:=
    \begin{bmatrix}
        Y_1 \\
        \vdots \\
        Y_n
    \end{bmatrix}
    \in\mathbb R^n.
\]
By \eqref{eq:gaussian-mixture-model}, the design matrix can be expressed as
\[
    X=\frac{Y\mu^\top}{\sqrt d}+\sigma U,
\]
where $U\in\mathbb R^{n\times d}$ has independent standard
Gaussian entries.
Define the residual vector $g^t\in\mathbb R^n$ by
\begin{align}
    g_i^t
    :=
    \frac{\Delta_i}{\rho}\ell'(Y_i,r_i^t)
    +
    \frac{\pi(1-\Delta_i)}{1-\rho}
        S_i^t\ell'(\widehat Y_i^t,r_i^t),
    \qquad i\in[\![1,n]\!].
    \label{eq:residual-map}
\end{align}
The recursion in Algorithm~\ref{alg:self-training} then becomes
\begin{align}
    \begin{cases}
    \displaystyle 
        r^t
        =
        b^t1_n+\frac{Xw^t}{\sqrt d},\\[5pt]
        \displaystyle 
        b^{t+1}
        =
        b^t-\eta\frac{1_n^\top g^t}{n},\\[5pt]
        \displaystyle 
        w^{t+1}
        =
        w^t-\eta\left[
            \lambda\nabla J(w^t)
            +\frac{\sqrt d}{n}X^\top g^t
        \right],
    \end{cases}
    \hspace{-10pt}
    \iff 
    \begin{cases}
    \displaystyle 
        r^t
        =
        b^t1_n+ m^t Y + \sigma q^t,\\[10pt]
        \displaystyle 
        b^{t+1}
        =
        b^t-\eta \zeta^t ,\\[5pt]
        \displaystyle 
        w^{t+1}
        =
        w^t-\eta\left[
            \lambda\nabla J(w^t)
            +
            \chi^t \mu 
            +\frac{\sigma}{\sqrt \delta} p^t
        \right],
    \end{cases}
    \label{eq:finite-dimensional-mean-field-iterations}
\end{align}
where
\begin{align}
    m^t := \frac{\mu^\top w^t}{d}, \quad q^t :=  \frac{U w^t}{\sqrt d}
    , \quad 
    \zeta^t := \frac{1_n^\top g^t}{n}, \quad \chi^t := \frac{Y^\top g^t}{n}, \quad p^t := \frac{U^\top g^t}{\sqrt n}.
    \label{eq:finite-dim-macro}
\end{align}
For a vector sequence $(a^s)_{0 \le s \le T}$ in $\mathbb R^m$,
we collect its history through time $t$ in the trajectory matrix
\[
    A_t:=
    \begin{bmatrix}
        a^0 & \cdots & a^t
    \end{bmatrix}
    \in\mathbb R^{m\times(t+1)},
    \qquad 0\leq t\leq T,
\]
with $A_{-1}$ understood to be empty.
We write $W_t$, $P_t$, $Q_t$, $G_t$, $R_t$ for the trajectory matrices
of $w^t$, $p^t$, $q^t$, $g^t$, $r^t$, respectively.
Row notation and the associated empirical measures are
defined in Appendix~\ref{subsec:empirical-measures}.

\begin{definition}[Algorithmic half-step filtrations]
\label{def:algorithmic-half-step-filtration}
For every \(t\in[\![0,T]\!]\), define
\begin{align}
    \mathcal F_t^{[w]}
    &:=
    \sigma\!\left(b^k,w^k:k\in[\![0,t]\!]\right),\\
    \mathcal F_t^{[g]}
    &:=
    \sigma\!\left(r^k:k\in[\![0,t]\!]\right).
\end{align}
Set $\mathcal F_{-1}^{[g]}:=\{\varnothing, \Omega\}$.
These are the canonical filtrations generated by the parameter trajectory and
the preactivation trajectory, respectively.
Since $b^0=b^{\mathrm{init}}$ and $w^0=w^{\mathrm{init}}$
are deterministic in the quenched model, \(\mathcal F_0^{[w]}\) is trivial.
We define the \emph{forward half-step filtration}
\((\mathcal F_t)_{t\in[\![0,T]\!]}\) and the \emph{backward half-step
filtration} \((\mathcal F_t^+)_{t\in[\![0,T]\!]}\) by
\begin{align}
    \mathcal F_t
    &:=
    \mathcal F_t^{[w]}\vee\mathcal F_{t-1}^{[g]},
    \label{eq:algorithmic-filtration-Ft}\\
    \mathcal F_t^+
    &:=
    \mathcal F_t^{[w]}\vee\mathcal F_t^{[g]}.
    \label{eq:algorithmic-filtration-Ft+}
\end{align}
The half-step filtrations satisfy
\[
    \mathcal F_0
    \subseteq
    \mathcal F_0^+
    \subseteq
    \mathcal F_1
    \subseteq
    \mathcal F_1^+
    \subseteq\cdots\subseteq
    \mathcal F_T
    \subseteq
    \mathcal F_T^+.
\]
The sigma-algebra \(\mathcal F_t\) represents the information available
immediately before the forward multiplication by \(U\), whereas
\(\mathcal F_t^+\) represents the information available immediately before
the backward multiplication by \(U^\top\) in \eqref{eq:finite-dimensional-mean-field-iterations}.
\end{definition}
The following lemma expresses the algorithmic filtrations in
terms of the forward and backward fluctuation histories.

\begin{lemma}[Equivalent fluctuation-history filtrations]
\label{lemma:equivalent-algorithmic-filtrations}
For every \(t\in[\![0,T]\!]\),
\begin{align}
    \mathcal F_t
    &=
    \sigma(Q_{t-1},P_{t-1})
    =
    \sigma(Q_{t-1},W_t,P_{t-1},G_{t-1}),
    \label{eq:equivalent-forward-filtration}\\
    \mathcal F_t^+
    &=
    \sigma(Q_t,P_{t-1})
    =
    \sigma(Q_t,W_t,P_{t-1},G_{t-1}).
    \label{eq:equivalent-backward-filtration}
\end{align}
In particular, $W_t,G_{t-1}$ are $\mathcal F_t$-measurable,
and $W_t,G_t$ are $\mathcal F_t^+$-measurable.
\end{lemma}

\begin{proof}
The initialization is deterministic, so $\mathcal F_0$ is trivial.
Given $\mathcal F_t$, the score equation shows that adjoining
$r^t$ is equivalent to adjoining $q^t$, since $\sigma>0$.
Similarly, given $\mathcal F_t^+$, the intercept update is
already determined, and adjoining $w^{t+1}$ is equivalent to
adjoining $p^t$, since $\eta\sigma>0$. Thus
\[
    \mathcal F_t^+
    =\mathcal F_t\vee\sigma(q^t),
    \qquad
    \mathcal F_{t+1}
    =\mathcal F_t^+\vee\sigma(p^t),
\]
with the second identity understood for $t<T$.
The claimed identities follow by induction, and the
measurability assertions follow from the definitions.
\end{proof}
For any matrix $A$, let $A^\dagger$ denote its
Moore--Penrose pseudoinverse and write
\[
    \Pi_A:=AA^\dagger,
    \qquad
    \Pi_A^\perp:=I-\Pi_A.
\]
Thus $\Pi_A$ is the orthogonal projector onto the column space
of $A$, with $\Pi_\emptyset=0$.
At each time $t$, define the normalized finite-dimensional Gram matrices
\[
 C_{t}^{[w]}:=\frac1dW_t^\top W_t,
 \qquad
 C_{t}^{[g]}:=\frac1nG_t^\top G_t.
\]
At the forward half-step, let
\begin{align}
 \alpha_{t}^{[w]}
 &:=
 (C_{t-1}^{[w]})^\dagger
 \frac{W_{t-1}^\top w^t}{d},
 &
 w_{\perp}^t
 &:=
 w^t-W_{t-1}\alpha_{t}^{[w]},
 \label{eq:finite-forward-projection-lemmas}\\
 \phi_{t}^{[w]}
 &:=
 (C_{t-1}^{[g]})^\dagger
 \frac{P_{t-1}^\top w_{\perp}^t}{d},
 &
 v_{t}^{[w]}
 &:=
 \frac{\|w_{\perp}^t\|}{\sqrt d}.
 \label{eq:finite-forward-memory-lemmas}
\end{align}
At \(t=0\), both coefficient vectors are empty and
\(w_{\perp}^0=w^0\). At the backward half-step, for \(t<T\), let
\begin{align}
 \alpha_{t}^{[g]}
 &:=
 (C_{t-1}^{[g]})^\dagger
 \frac{G_{t-1}^\top g^t}{n},
 &
 g_{\perp}^t
 &:=
 g^t-G_{t-1}\alpha_{t}^{[g]},
 \label{eq:finite-backward-projection-lemmas}\\
 \phi_{t}^{[g]}
 &:=
 (C_{t}^{[w]})^\dagger
 \frac{Q_t^\top g_{\perp}^t}{n},
 &
 v_{t}^{[g]}
 &:=
 \frac{\|g_{\perp}^t\|}{\sqrt n}.
 \label{eq:finite-backward-memory-lemmas}
\end{align}
At \(t=0\), \(\alpha_{0}^{[g]}\) is empty and
\(g_{\perp}^0=g^0\).

The notation above parallels that of
Definition~\ref{def:mean-field-recursion}.
Effective random quantities generally carry tildes and
deterministic mean-field quantities carry bars.
Gaussian innovations are an exception: $z_t^{[g]}$ and
$z_t^{[w]}$ denote finite-dimensional vectors in conditional
finite-dimensional representations and scalar standard
Gaussians in effective-process formulas.
Their dimensions and roles are specified at each use.
We use the empirical measure notation introduced in
Definition~\ref{definition:empirical-measures}.

For the effective trajectory vectors, write
\[
\begin{aligned}
    \tilde W_t&:=(\tilde w^0,\ldots,\tilde w^t)^\top,
    &\tilde P_t&:=(\tilde p^0,\ldots,\tilde p^t)^\top,\\
    \tilde Q_t&:=(\tilde q^0,\ldots,\tilde q^t)^\top,
    &\tilde G_t&:=(\tilde g^0,\ldots,\tilde g^t)^\top,\\
    \tilde R_t&:=(\tilde r^0,\ldots,\tilde r^t)^\top.
\end{aligned}
\]
Vectors indexed by $-1$ are empty.

For random elements $\xi,\xi'$ defined on a common extension
and a sub-$\sigma$-field $\mathcal G$, we write
$\xi\stackrel d=_{\mathcal G}\xi'$ when
\[
    \mathscr L(\xi\mid\mathcal G)
    =
    \mathscr L(\xi'\mid\mathcal G)
    \qquad\text{almost surely}.
\]
The conditioning $\sigma$-field is identified with its pullback
to the extension. In particular,
$(A,\xi)\stackrel d=(A,\xi')$ for every
$\mathcal G$-measurable random element $A$.
Auxiliary Gaussian variables are constructed on product
extensions independently of the corresponding conditioning
history.

\begin{proposition}[Master recursion and conditional Gaussian decomposition]
\label{proposition:master-recursion}
Suppose that
Assumption~\ref{assumption:asymptotic-regime},
Assumption~\ref{assumption:quenched-initialization},
and Assumption~\ref{assumption:regularity-of-loss-and-penalty} hold.

The well-definedness, moment, and nondegeneracy conclusions
of Lemma~\ref{lemma:mf-proof-nondegeneracy} hold.
In each conditional-law representation below, $\check U$ denotes
an auxiliary standard Gaussian matrix independent of the
corresponding conditioning sigma-algebra.
The following assertions hold.

\begin{description}[wide]
    \item[$(\mathsf F_t)$ Forward half-step:]
    For $0\leq t\leq T$:
    \begin{enumerate}[label=(\roman*),wide]
        \item \emph{Information and conditional matrix law.}
        The matrices $W_t,G_{t-1}$ are $\mathcal F_t$-measurable by
        Lemma~\ref{lemma:equivalent-algorithmic-filtrations}.
        The conditional matrix law and its mean are given by
        \eqref{eq:mf-proof-matrix-forward}
        and \eqref{eq:mf-proof-mean-forward}.

        \item \emph{Projection coefficients and nondegeneracy.}
        As $d\to\infty$,
        \begin{align}
            v_t^{[w]}
            &\xrightarrow{\mathbb P}\bar v_t^{[w]}>0,
            &
            C_t^{[w]}
            &\xrightarrow{\mathbb P}\bar C_t^{[w]}\succ0,
            \label{eq:master-forward-nondegeneracy}
            \\
            \alpha_t^{[w]}
            &\xrightarrow{\mathbb P}\bar\alpha_t^{[w]},
            &
            \phi_t^{[w]}
            &\xrightarrow{\mathbb P}\bar\phi_t^{[w]}.
            \label{eq:master-forward-coefficient-convergence}
        \end{align}

        \item \emph{Conditional forward fluctuation.}
        One has
        \begin{align}
            q^t
            &\stackrel d=_{\mathcal F_t}
            Q_{t-1}\alpha_t^{[w]}
            +\frac1{\sqrt\delta}G_{t-1}\phi_t^{[w]}
            +v_t^{[w]}z_t^{[g]}
            +\varepsilon_t^{[g]},
            \label{eq:conditional-q-decomposition-lemmas}
        \end{align}
        where $z_t^{[g]}\sim\mathcal N(0,I_n)$ is independent
        of $\mathcal F_t$ and
        \(
            \varepsilon_t^{[g]}
            :=
            -v_t^{[w]}\Pi_{G_{t-1}}z_t^{[g]}
        \).
        Moreover,
        \begin{equation}
            \mathbb E\left[
                \frac{\|\varepsilon_t^{[g]}\|^2}{n}
                \Bigm|\mathcal F_t
            \right]
            \leq
            (v_t^{[w]})^2\frac tn,
            \label{eq:forward-projection-error-lemmas}
        \end{equation}
        and, by tightness of $v_t^{[w]}$, we have
        \(
            \|\varepsilon_t^{[g]}\|^2 / n
            \xrightarrow{\mathbb P}0
        \).

        \item \emph{Sample-side empirical extension.}
        The joint empirical law satisfies
        \begin{equation}
            W_2 \left(
            \nu(Y,\Delta,Q_t,G_{t-1},R_t)
            , \,  
            \mathscr L(
                \tilde Y,\tilde\Delta,
                \tilde Q_t,\tilde G_{t-1},\tilde R_t
                )
            \right)
            \xrightarrow{\mathbb P}
            0.
            \label{eq:master-forward-empirical-law}
        \end{equation}
        In addition,
        \(
            m^t\xrightarrow{\mathbb P}\bar m^t
        \).
        For $t<T$, one also has
        \begin{equation}
            W_2 \left(
            \nu{(Y,\Delta,Q_t,G_t,R_t)}
            , \,
            \mathscr L(
                \tilde Y,\tilde\Delta,
                \tilde Q_t,\tilde G_t,\tilde R_t
            )
            \right)
            \xrightarrow{\mathbb P}
            0
            ,
            \label{eq:master-forward-residual-extension}
        \end{equation}
        and
        \begin{equation}
            \zeta^t\xrightarrow{\mathbb P}\bar\zeta^t,
            \qquad
            \chi^t\xrightarrow{\mathbb P}\bar\chi^t,
            \qquad
            b^{t+1}\xrightarrow{\mathbb P}\bar b^{t+1}.
            \label{eq:master-backward-macroscopic-convergence}
        \end{equation}
    \end{enumerate}

    \item[$(\mathsf B_t)$ Backward half-step:]
    For $0\leq t<T$:
    \begin{enumerate}[label=(\roman*),wide]
        \item \emph{Information and conditional matrix law.}
        The matrices $W_t,G_t$ are $\mathcal F_t^+$-measurable by
        Lemma~\ref{lemma:equivalent-algorithmic-filtrations}.
        The conditional matrix law and its mean are given by
        \eqref{eq:mf-proof-matrix-backward}
        and \eqref{eq:mf-proof-mean-backward}.

        \item \emph{Projection coefficients and nondegeneracy.}
        As $d\to\infty$,
        \begin{align}
            v_t^{[g]}
            &\xrightarrow{\mathbb P}\bar v_t^{[g]}>0,
            &
            C_t^{[g]}
            &\xrightarrow{\mathbb P}\bar C_t^{[g]}\succ0,
            \label{eq:master-residual-nondegeneracy}
            \\
            \alpha_t^{[g]}
            &\xrightarrow{\mathbb P}\bar\alpha_t^{[g]},
            &
            \phi_t^{[g]}
            &\xrightarrow{\mathbb P}\bar\phi_t^{[g]}.
            \label{eq:master-backward-coefficient-convergence}
        \end{align}

        \item \emph{Conditional backward fluctuation.}
        One has
        \begin{align}
            p^t
            &\stackrel d=_{\mathcal F_t^+}
            P_{t-1}\alpha_t^{[g]}
            +\sqrt\delta\,W_t\phi_t^{[g]}
            +v_t^{[g]}z_t^{[w]}
            +\varepsilon_t^{[w]},
            \label{eq:conditional-p-decomposition-lemmas}
        \end{align}
        where $z_t^{[w]}\sim\mathcal N(0,I_d)$ is independent
        of $\mathcal F_t^+$ and
        \(
            \varepsilon_t^{[w]}
            :=
            -v_t^{[g]}\Pi_{W_t}z_t^{[w]}.
        \)
        Moreover,
        \begin{equation}
            \mathbb E\left[
                \frac{\|\varepsilon_t^{[w]}\|^2}{d}
                \Bigm|\mathcal F_t^+
            \right]
            \leq
            (v_t^{[g]})^2\frac{t+1}{d},
            \label{eq:backward-projection-error-lemmas}
        \end{equation}
        and, by tightness of $v_t^{[g]}$, we have
        \(
            \|\varepsilon_t^{[w]}\|^2 / d
            \xrightarrow{\mathbb P}0
        \).

        \item \emph{Parameter-side empirical extension.}
        The joint empirical law satisfies
        \begin{equation}
            W_2 \left( 
            \nu{(\mu,W_{t+1},P_t)} 
            , \, 
            \mathscr L(
                \tilde\mu,\tilde W_{t+1},\tilde P_t
            )
            \right)
            \xrightarrow{\mathbb P}
            0
            .
            \label{eq:master-backward-empirical-law}
        \end{equation}
    \end{enumerate}
\end{description}
\end{proposition}

\subsection{Supporting lemmas}
\label{subsec:supporting-lem}

    \subsubsection{Empirical convergence}

        We establish general lemmas,
        Lemmas~\ref{lemma:mf-proof-empirical-stability}
        and~\ref{lemma:mf-proof-gaussian-extension},
        used in the proof of the master recursion in 
        Proposition~\ref{proposition:master-recursion}.
        Throughout this subsection, the superscript $(m)$
        denotes the number of rows, with convergence understood
        along the underlying sequence of arrays.
        We apply these lemmas with $m=d$ for parameter-side
        empirical laws and $m=n(d)$ for sample-side empirical laws.

        All convergence-in-probability statements below have
        deterministic limits and depend only on the joint law of
        the objects at each dimension. We may therefore realize
        the objects at all dimensions on a common product
        probability space, preserving their within-dimension joint
        laws. We use this convention when extracting almost-surely
        convergent subsequences.

        \begin{lemma}[Stability of empirical laws]
        \label{lemma:mf-proof-empirical-stability}
        Fix an integer $k\geq1$ and a closed set
        $\mathcal X\subseteq\mathbb R^k$.
        Consider a sequence of random matrices, with row counts
        $m\to\infty$, of the form
        \[
            H^{(m)}
            :=
            \begin{bmatrix}
                (h_1^{(m)})^\top\\
                \vdots\\
                (h_m^{(m)})^\top
            \end{bmatrix}
            \in\mathbb R^{m\times k},
            \qquad
            h_i^{(m)}\in\mathcal X
            \quad\text{almost surely}.
        \]
        Suppose
        \[
            W_2\left(
                \nu(H^{(m)}),
                \mathscr L(\tilde H)
            \right)
            \xrightarrow{\mathbb P}0,
            \qquad
            \mathscr L(\tilde H)\in\mathscr P_2(\mathcal X).
        \]
        Then the following statements hold.
        
        \begin{enumerate}[wide]
            \item \emph{Moment and Gram-matrix convergence.}
            One has
            \[
                \frac1m\sum_{i=1}^m h_i^{(m)}
                \xrightarrow{\mathbb P}
                \tilde{\mathbb E}[\tilde H],
                \qquad
                K^{(m)}
                :=
                \frac1m(H^{(m)})^\top H^{(m)}
                \xrightarrow{\mathbb P}
                K
                :=
                \tilde{\mathbb E}[\tilde H\tilde H^\top].
            \]
            In particular,
            \[
                K\succ0
                \quad\Longrightarrow\quad
                (K^{(m)})^\dagger
                \xrightarrow{\mathbb P}K^{-1}.
            \]
        
            \item \emph{Stability under rowwise transformations.}
            Let $\beta^{(m)}\in\mathbb R^p$ be random vectors satisfying
            \[
                \beta^{(m)}
                \xrightarrow{\mathbb P}
                \bar\beta\in\mathbb R^p,
            \]
            where $p$ and $\bar\beta$ are fixed.
            Let
            \(
                F:\mathbb R^p\times\mathcal X\longrightarrow\mathbb R^q
            \)
            be measurable, with $q$ fixed.
            Suppose $F$ is continuous at $(\bar\beta,x)$ for
            $\mathscr L(\tilde H)$-almost every $x$, with continuity
            understood relative to $\mathbb R^p\times\mathcal X$, and
            \[
                \exists\,C,\varepsilon>0
                \quad
                \forall\,\|\beta-\bar\beta\|<\varepsilon,
                \quad
                \forall\,x\in\mathcal X,
                \quad
                \|F(\beta,x)\|\leq C(1+\|x\|).
            \]
            Define $F^{(m)}\in\mathbb R^{m\times q}$ by
            \[
                F^{(m)}(i)
                :=
                F(\beta^{(m)},h_i^{(m)})^\top.
            \]
            Then
            \[
                W_2\left(
                    \nu(H^{(m)},F^{(m)}),
                    \mathscr L\bigl(
                        \tilde H,F(\bar\beta,\tilde H)
                    \bigr)
                \right)
                \xrightarrow{\mathbb P}0.
            \]
        \end{enumerate}
        \end{lemma}
        
        \begin{proof}
        Recall that convergence in probability is equivalent to the
        following subsequence property: every subsequence admits a
        further subsequence converging almost surely to the same limit
        \cite[Lemma~5.2]{kallenberg2021foundations}.
        
        Starting from any subsequence $(m_k)_k$, extract a further
        subsequence $(m_{k_j})_j$ along which the assumed Wasserstein convergence
        and, for part~2, coefficient convergence hold almost surely.
        Fix an outcome $\omega \in \Omega$ where these convergences hold and all the rows
        belong to $\mathcal X$.
        Suppressing the subsequence indices by denoting $m := m_{k_j}$, the realized matrices
        and coefficients form deterministic sequences satisfying
        \[
            W_2\bigl(\nu(H^{(m)}),\mathscr L(\tilde H)\bigr)
            \longrightarrow0,
            \qquad
            \beta^{(m)}\longrightarrow\bar\beta
            .
        \]
        We prove the conclusions for these deterministic sequences.
        By the coupling characterization of Wasserstein convergence
        (see \cite[Definition~6.8 and Theorem~6.9]{villani2009optimal},
        together with
        \cite[Theorems~5.12 and~5.31 and Lemma~5.11(ii)]
        {kallenberg2021foundations}),
        there exist random vectors $X^{(m)},X$ on an auxiliary
        probability space such that
        \[
            \mathscr L(X^{(m)})=\nu(H^{(m)}),
            \qquad
            \mathscr L(X)=\mathscr L(\tilde H),
            \qquad
            \mathbb E\|X^{(m)}-X\|^2\longrightarrow0.
        \]
        Their marginal laws are supported on $\mathcal X$.
        Consequently, $X$ and all the $X^{(m)}$ belong to
        $\mathcal X$ simultaneously almost surely.
        
        \begin{enumerate}[wide]
            \item \emph{Moment and Gram-matrix convergence.}
            Since
            \[
                \mathbb E[X^{(m)}]
                =
                \frac1m\sum_{i=1}^m h_i^{(m)},
                \qquad
                \mathbb E[X]
                =
                \tilde{\mathbb E}[\tilde H],
            \]
            the first-moment assertion follows from
            \[
                \|\mathbb E[X^{(m)}]-\mathbb E[X]\|
                \leq
                \left(
                    \mathbb E\|X^{(m)}-X\|^2
                \right)^{1/2}
                \longrightarrow0.
            \]
            For the second moments, write
            \[
                X^{(m)}(X^{(m)})^\top-XX^\top
                =
                (X^{(m)}-X)(X^{(m)})^\top
                +
                X(X^{(m)}-X)^\top.
            \]
            Taking expectations and applying Cauchy--Schwarz yields
            \[
            \begin{aligned}
                \left\|
                    \mathbb E[X^{(m)}(X^{(m)})^\top]
                    -
                    \mathbb E[XX^\top]
                \right\|_{\mathrm F}
                \leq
                \left(
                    \mathbb E\|X^{(m)}-X\|^2
                \right)^{1/2}
                \left[
                    \left(\mathbb E\|X^{(m)}\|^2\right)^{1/2}
                    +
                    \left(\mathbb E\|X\|^2\right)^{1/2}
                \right]
                \longrightarrow0.
            \end{aligned}
            \]
            Here the bracketed factor is bounded because
            $X^{(m)}\to X$ in $L^2$.
            The two moment matrices are $K^{(m)}$ and $K$,
            respectively, proving Gram-matrix convergence.
            If $K\succ0$, then, for all sufficiently large $m$,
            \[
                \lambda_{\min}(K^{(m)})
                \geq
                \lambda_{\min}(K)
                -
                \|K^{(m)}-K\|_{\mathrm{op}}
                >
                \frac{\lambda_{\min}(K)}2.
            \]
            Thus $(K^{(m)})^\dagger=(K^{(m)})^{-1}$ eventually,
            and continuity of matrix inversion gives
            \[
                (K^{(m)})^\dagger\longrightarrow K^{-1}.
            \]
            Applying the same argument to coordinate projections
            proves the assertions for fixed subcollections of columns.
        
            \item \emph{Stability under rowwise transformations.}
            The deterministic coefficient convergence and the coupling
            imply
            \[
                (\beta^{(m)},X^{(m)})
                \xrightarrow{\mathbb P}
                (\bar\beta,X)
            \]
            as random elements of $\mathbb R^p\times\mathcal X$,
            equipped with its relative Euclidean metric.
            Since $F$ is continuous at $(\bar\beta,X)$ almost surely,
            the continuous mapping theorem
            \cite[Lemma~5.3]{kallenberg2021foundations} gives
            \[
                F(\beta^{(m)},X^{(m)})
                \xrightarrow{\mathbb P}
                F(\bar\beta,X).
            \]        
            To strengthen this to $L^2$ convergence, observe that
            $(\|X^{(m)}\|^2)_m$ is uniformly integrable because
            $X^{(m)}\to X$ in $L^2$.
            For all sufficiently large $m$,
            $\|\beta^{(m)}-\bar\beta\|<\varepsilon$.
            Since the coupled variables belong to $\mathcal X$,
            the growth assumption gives
            \[
            \begin{aligned}
                \|F(\beta^{(m)},X^{(m)})
                    -F(\bar\beta,X)\|^2
                \leq
                2\|F(\beta^{(m)},X^{(m)})\|^2
                +
                2\|F(\bar\beta,X)\|^2
                \leq
                4C^2\bigl(
                    2+\|X^{(m)}\|^2+\|X\|^2
                \bigr).
            \end{aligned}
            \]
            Thus the squared differences are uniformly integrable
            for all sufficiently large $m$.
            Their convergence to zero in probability and Vitali's
            convergence theorem
            \cite[Theorem~5.12]{kallenberg2021foundations}
            yield
            \[
                \mathbb E\left[
                    \|F(\beta^{(m)},X^{(m)})
                        -F(\bar\beta,X)\|^2
                \right]
                \longrightarrow0.
            \]
            The joint vectors
            \[
                (X^{(m)},F(\beta^{(m)},X^{(m)}))
                \qquad\text{and}\qquad
                (X,F(\bar\beta,X))
            \]
            have laws $\nu(H^{(m)},F^{(m)})$ and
            $\mathscr L(\tilde H,F(\bar\beta,\tilde H))$,
            respectively.
            Their coupling therefore gives
            \[
            \begin{aligned}
                W_2^2\left(
                    \nu(H^{(m)},F^{(m)}),
                    \mathscr L(\tilde H,F(\bar\beta,\tilde H))
                \right)
                \leq
                \mathbb E\|X^{(m)}-X\|^2
                +
                \mathbb E\left[
                    \|F(\beta^{(m)},X^{(m)})
                        -F(\bar\beta,X)\|^2
                \right]
                \longrightarrow0.
            \end{aligned}
            \]
        \end{enumerate}
        Applying these deterministic conclusions along the extracted
        subsequences gives almost-sure convergence there.
        Since the initial subsequence was arbitrary, the subsequence
        characterization proves both assertions in probability.
        \end{proof}

        \begin{lemma}[Conditional Gaussian extension]
        \label{lemma:mf-proof-gaussian-extension}
        Fix an integer $k\geq1$.
        Along a sequence of row counts $m\to\infty$, let
        \[
            H^{(m)}\in\mathbb R^{m\times k},
            \qquad
            c^{(m)}\in\mathbb R^k,
            \qquad
            v^{(m)}\geq0
        \]
        be $\mathcal G^{(m)}$-measurable, and let
        $x^{(m)}\in\mathbb R^m$ be random.
        Suppose, as $m\to\infty$,
        \[
            W_2\left(
                \nu(H^{(m)}),
                \mathscr L(\tilde H)
            \right)\xrightarrow{\mathbb P}0,
            \qquad
            c^{(m)}\xrightarrow{\mathbb P}\bar c,
            \qquad
            v^{(m)}\xrightarrow{\mathbb P}\bar v,
        \]
        where
        $\mathscr L(\tilde H)\in\mathscr P_2(\mathbb R^k)$,
        and $\bar c\in\mathbb R^k$, $\bar v\geq0$ are deterministic.
        Assume that, for each $m$,
        \begin{align}
            x^{(m)}
            \stackrel d=_{\mathcal G^{(m)}}
            H^{(m)}c^{(m)}+v^{(m)}z^{(m)}+e^{(m)},
            \label{eq:conditional-representation}
        \end{align}
        where $z^{(m)}\sim\mathcal N(0,I_m)$ is independent of
        $\mathcal G^{(m)}$ and $e^{(m)}\in\mathbb R^m$ satisfies
        \(
            \|e^{(m)}\|^2 / m
            \xrightarrow{\mathbb P}0
        \).
        The conditional-law notation is as defined before Proposition~\ref{proposition:master-recursion}.
        Then,
        \[
            W_2\left(
                \nu(H^{(m)},x^{(m)}),
                \mathscr L\bigl(
                    \tilde H,\tilde H^\top\bar c+\bar v Z
                \bigr)
            \right)
            \xrightarrow{\mathbb P}0,
        \]
        where $Z\sim\mathcal N(0,1)$ is independent of $\tilde H$.
        \end{lemma}

        \begin{proof}
        Work on a probability space supporting the right-hand side of
        the conditional representation \eqref{eq:conditional-representation}.
        Write $(h_i^{(m)})^\top$ for the $i$-th row of $H^{(m)}$ and set
        \[
            y^{(m)}
            :=
            H^{(m)}c^{(m)}+v^{(m)}z^{(m)}.
        \]
        Let $Z\sim\mathcal N(0,1)$ be independent of $\tilde H$.
        
        \begin{enumerate}[wide]
            \item \emph{Adjoining the Gaussian column.}
            We first establish
            \[
                W_2\left(
                    \nu(H^{(m)},z^{(m)}),
                    \mathscr L(\tilde H,Z)
                \right)
                \xrightarrow{\mathbb P}0.
            \]
            Fix a bounded Lipschitz function
            $f:\mathbb R^{k+1}\to\mathbb R$ and define
            \[
                A^{(m)}(f)
                :=
                \frac1m\sum_{i=1}^m
                    f(h_i^{(m)},z_i^{(m)}),
                \qquad
                f_\gamma(h):=\mathbb E_Z[f(h,Z)].
            \]
            The function $f_\gamma$ is bounded and Lipschitz, since
            \[
                \|f_\gamma\|_\infty\leq\|f\|_\infty,
                \qquad
                |f_\gamma(h)-f_\gamma(h')|
                \leq\operatorname{Lip}(f)\|h-h'\|.
            \]
            Conditionally on $\mathcal G^{(m)}$, the rows
            $h_i^{(m)}$ are fixed and the coordinates $z_i^{(m)}$
            are independent standard Gaussian variables. Therefore,
            \[
                \mathbb E[A^{(m)}(f)\mid\mathcal G^{(m)}]
                =
                \frac1m\sum_{i=1}^m f_\gamma(h_i^{(m)})
                =
                \int f_\gamma\,d\nu(H^{(m)}).
            \]
            The Lipschitz bound gives
            \[
            \begin{aligned}
                \left|
                    \mathbb E[A^{(m)}(f)\mid\mathcal G^{(m)}]
                    -\tilde{\mathbb E}[f(\tilde H,Z)]
                \right|
                =
                \left|
                    \int f_\gamma\,d\nu(H^{(m)})
                    -\tilde{\mathbb E}[f_\gamma(\tilde H)]
                \right|
                \leq
                \operatorname{Lip}(f)\,
                W_2\bigl(\nu(H^{(m)}),\mathscr L(\tilde H)\bigr)
                \xrightarrow{\mathbb P}0.
            \end{aligned}
            \]
            Here independence of $Z$ and $\tilde H$ identifies
            $\tilde{\mathbb E}[f_\gamma(\tilde H)]$ with
            $\tilde{\mathbb E}[f(\tilde H,Z)]$.
            Conditional independence also yields
            \[
            \begin{aligned}
                \operatorname{Var}
                \bigl(A^{(m)}(f)\mid\mathcal G^{(m)}\bigr)
                =
                \frac1{m^2}\sum_{i=1}^m
                \operatorname{Var}
                \bigl(f(h_i^{(m)},z_i^{(m)})
                      \mid\mathcal G^{(m)}\bigr)
                \leq\frac{\|f\|_\infty^2}{m}.
            \end{aligned}
            \]
            Consequently, for every $\varepsilon>0$,
            \[
                \mathbb P\left(
                    \left|
                        A^{(m)}(f)
                        -\mathbb E[A^{(m)}(f)\mid\mathcal G^{(m)}]
                    \right|>\varepsilon
                \right)
                \leq\frac{\|f\|_\infty^2}{m\varepsilon^2}.
            \]
            Combining the conditional mean and variance estimates gives
            \[
                A^{(m)}(f)
                \xrightarrow{\mathbb P}
                \tilde{\mathbb E}[f(\tilde H,Z)].
            \]
            The second moment of the joint empirical measure is
            \[
                \int\|(h,z)\|^2\,\nu(H^{(m)},z^{(m)})(dh,dz)
                =
                \frac{\|H^{(m)}\|_{\mathrm F}^2}{m}
                +\frac{\|z^{(m)}\|^2}{m}.
            \]
            By part~1 of
            Lemma~\ref{lemma:mf-proof-empirical-stability},
            \[
                \frac{\|H^{(m)}\|_{\mathrm F}^2}{m}
                \xrightarrow{\mathbb P}
                \tilde{\mathbb E}\|\tilde H\|^2.
            \]
            Moreover,
            \[
                \mathbb E\left[\frac{\|z^{(m)}\|^2}{m}\right]=1,
                \qquad
                \operatorname{Var}\left(
                    \frac{\|z^{(m)}\|^2}{m}
                \right)=\frac2m,
            \]
            so $\|z^{(m)}\|^2/m\xrightarrow{\mathbb P}1$.
            Thus the joint empirical second moment converges in
            probability to $\tilde{\mathbb E}\|\tilde H\|^2+1$.
        
            To conclude, the preceding estimates establish convergence
            in probability of the empirical integrals against every
            bounded Lipschitz test function, together with convergence
            of the second moments.
            Proposition~\ref{proposition:wasserstein-convergence-characterization}
            therefore yields the claimed convergence in $W_2$.
        
            \item \emph{Applying the affine transformation.}
            Apply part~2 of
            Lemma~\ref{lemma:mf-proof-empirical-stability} to
            \[
                [\,H^{(m)}\mid z^{(m)}\,],
                \qquad
                \beta^{(m)}:=(c^{(m)},v^{(m)}),
            \]
            and the map
            \(
                F((c,v),(h,z)):=h^\top c+vz
            \).
            The coefficients converge in probability to
            $\bar\beta:=(\bar c,\bar v)$.
            The map $F$ is continuous and, whenever
            $\|\beta-\bar\beta\|<1$,
            \[
                |F(\beta,(h,z))|
                \leq\|\beta\|\,\|(h,z)\|
                \leq(\|\bar\beta\|+1)\|(h,z)\|.
            \]
            Hence it satisfies the required uniform linear-growth
            condition near $\bar\beta$.
            The lemma gives
            \[
                W_2\left(
                    \nu(H^{(m)},z^{(m)},y^{(m)}),
                    \mathscr L\bigl(
                        \tilde H,Z,\tilde H^\top\bar c+\bar v Z
                    \bigr)
                \right)
                \xrightarrow{\mathbb P}0.
            \]
            By Lemma~\ref{lemma:wasserstein-lipschitz-pushforward},
            the coordinate projection $(h,z,y)\mapsto(h,y)$ cannot
            increase $W_2$.
            It follows that
            \[
                W_2\left(
                    \nu(H^{(m)},y^{(m)}),
                    \mu_\infty
                \right)
                \xrightarrow{\mathbb P}0,
                \qquad
                \mu_\infty
                :=
                \mathscr L(\tilde H,\tilde H^\top\bar c+\bar v Z).
            \]
        
            \item \emph{Including the negligible error.}
            For each realization, couple rows with the same index:
            \[
                \Gamma^{(m)}
                :=
                \frac1m\sum_{i=1}^m
                \delta_{\bigl(
                    (h_i^{(m)}, \, y_i^{(m)}+e_i^{(m)}), \,
                    (h_i^{(m)}, \, y_i^{(m)})
                \bigr)}.
            \]
            The measure $\Gamma^{(m)}$ is a coupling of
            $\nu(H^{(m)}, \,y^{(m)}+e^{(m)})$ and
            $\nu(H^{(m)}, \,y^{(m)})$.
            The history coordinates coincide in each coupled pair,
            so its quadratic transportation cost is
            \[
                \frac1m\sum_{i=1}^m |e_i^{(m)}|^2
                =
                \frac{\|e^{(m)}\|^2}{m}.
            \]
            By the definition of $W_2$,
            \[
                W_2^2\left(
                    \nu(H^{(m)},y^{(m)}+e^{(m)}),
                    \nu(H^{(m)},y^{(m)})
                \right)
                \leq\frac{\|e^{(m)}\|^2}{m}
                \xrightarrow{\mathbb P}0.
            \]
            The triangle inequality and the preceding step yield
            \[
                W_2\left(
                    \nu(H^{(m)},y^{(m)}+e^{(m)}),
                    \mu_\infty
                \right)
                \xrightarrow{\mathbb P}0.
            \]
            This argument is pointwise in the realized arrays and
            requires no independence between $e^{(m)}$ and
            $z^{(m)}$.
        
            \item \emph{Transferring the conclusion to $x^{(m)}$.}
            Since $H^{(m)}$ is $\mathcal G^{(m)}$-measurable,
            equality of the conditional laws of $x^{(m)}$ and
            $y^{(m)}+e^{(m)}$ also identifies their joint laws
            with $H^{(m)}$.
            Indeed, for every bounded Borel function
            $\varphi:\mathbb R^{m\times k}\times\mathbb R^m
            \to\mathbb R$,
            \[
            \begin{aligned}
                \mathbb E[\varphi(H^{(m)},x^{(m)})]
                &=
                \mathbb E\left[
                    \mathbb E[
                        \varphi(H^{(m)},x^{(m)})
                        \mid\mathcal G^{(m)}
                    ]
                \right]
                =
                \mathbb E\left[
                    \mathbb E[
                        \varphi(H^{(m)},y^{(m)}+e^{(m)})
                        \mid\mathcal G^{(m)}
                    ]
                \right]\\
                &=
                \mathbb E[
                    \varphi(H^{(m)},y^{(m)}+e^{(m)})
                ].
            \end{aligned}
            \]
            Therefore
            \[
                (H^{(m)},x^{(m)})
                \stackrel d=
                (H^{(m)},y^{(m)}+e^{(m)}).
            \]
            Since $\mu_\infty$ is deterministic, their empirical
            Wasserstein distances to $\mu_\infty$ have the same
            distribution. In particular, for every $\varepsilon>0$,
            \[
            \begin{aligned}
                &\mathbb P\left(
                    W_2\bigl(\nu(H^{(m)},x^{(m)}),\mu_\infty\bigr)
                    >\varepsilon
                \right)
                =
                \mathbb P\left(
                    W_2\bigl(
                        \nu(H^{(m)},y^{(m)}+e^{(m)}),\mu_\infty
                    \bigr)>\varepsilon
                \right)
                \longrightarrow0.
            \end{aligned}
            \]
            This is the desired conclusion.
        \end{enumerate}
        \end{proof}

    \subsubsection{Gaussian conditioning}

    We use the Gaussian conditioning approach employed in
    \cite{bolthausen2014} and
    \cite[Section~III-H, Lemmas~10--11]{bayati2011dynamics}.
    Our linear-conditioning lemma also permits rank-deficient histories.
    The key idea is to separate the design matrix into its
    conditional mean, determined by the observations revealed so far,
    and a Gaussian fluctuation in the remaining unconstrained directions.
    We first state
    Lemma~\ref{lemma:mf-proof-linear-conditioning}
    in a finite-dimensional Euclidean space.
    We then apply it recursively to $\mathbb R^{n\times d}$
    with its Frobenius inner product, obtaining the conditional
    laws in Proposition~\ref{proposition:mf-proof-conditional-matrix}.

        \begin{lemma}[Gaussian conditioning on a linear observation]
        \label{lemma:mf-proof-linear-conditioning}
        Let $E,E'$ be finite-dimensional Euclidean spaces,
        let $F\subseteq E$ be a linear subspace, and let
        \[
            X=m+\pi_F (Z),
            \qquad
            Z\sim\mathcal N(0,I_E),
        \]
        where $m\in E$ is deterministic and $\pi_F$ is the
        orthogonal projection onto $F$.
        Let $T:E\to E'$ be a deterministic linear map and set
        \[
            K:=F\cap\ker T.
        \]
        For $\mathscr L(T(X)) $-almost every $y$, the conditional
        law of $X$ given $T(X)=y$ is
        \[
            \mathcal N(m_y,\pi_K),
        \]
        where $\pi_K$ is the orthogonal projection onto $K$
        and $m_y$ is uniquely determined by
        \[
            m_y\in m+F,
            \qquad
            T(m_y)=y,
            \qquad
            m_y-m\perp K.
        \]
        If $m\perp F$, then $m_y\perp K$.
        \end{lemma}

        \begin{proof}
        Set
        \(
            L:=F\cap K^\perp
        \).
        We proceed in three steps.
        
        \begin{enumerate}[wide]
            \item \emph{Orthogonal Gaussian decomposition.}
            We have the decomposition:
            \[
                F=K\oplus L,
                \qquad
                \pi_F=\pi_K+\pi_L.
            \]
            Define $\xi:=\pi_K Z$, and $\eta:=\pi_L Z$.
            Then
            \[
                X=m+\xi+\eta.
            \]
            The vectors $\xi$ and $\eta$ are jointly Gaussian as the projections are linear and $Z$ is Gaussian.
            They are centered, and their covariance operators satisfy
            \[
                \operatorname{Var}(\xi)=\pi_K,
                \qquad
                \operatorname{Var}(\eta)=\pi_L,
                \qquad
                \operatorname{Cov}(\xi,\eta)
                =\pi_K\pi_L=0.
            \]
            Consequently,
            \[
                \xi\sim\mathcal N(0,\pi_K),
                \qquad
                \eta\sim\mathcal N(0,\pi_L),
                \qquad
                \xi\independent\eta,
            \]
            by the characterization of Gaussian laws by their means
            and covariances and the independence criterion for
            jointly Gaussian vectors
            \cite[Lemma~14.1]{kallenberg2021foundations}.
        
            \item \emph{The observed component and the conditional mean.}
            Consider the linear map $A :=T|_L$, where
            \[
                T|_L:L\longrightarrow T(F).
            \]
            This map is injective, because
            \[
                \ker A=L\cap\ker T=L\cap K=\{0\}.
            \]
            It is also surjective: if $f=k+l$ with $k\in K$
            and $l\in L$, then
            \[
                Tf=Tl=Al.
            \]
            Hence $A$ is a linear bijection between finite-dimensional
            vector spaces, and its inverse is continuous.
            Write $Y:=T(X)$. Since $\xi\in K\subseteq\ker T$,
            \[
                Y=Tm+A\eta.
            \]
            In particular,
            \[
                Y\in Tm+T(F),
                \qquad
                \eta=A^{-1}(Y-Tm).
            \]
            For every $y\in Tm+T(F)$, define
            \[
                m_y:=m+A^{-1}(y-Tm).
            \]
            By construction,
            \[
                m_y\in m+F,
                \qquad
                Tm_y=y,
                \qquad
                m_y-m\in L\subseteq K^\perp.
            \]
            These conditions determine $m_y$ uniquely.
            Indeed, if $x\in m+F$ also satisfies $Tx=y$ and
            $x-m\perp K$, then
            \[
                x-m_y\in F\cap\ker T=K.
            \]
            Moreover, both $x-m$ and $m_y-m$ belong to $K^\perp$,
            so $x-m_y\in K^\perp$. Therefore $x-m_y=0$.
        
            \item \emph{Identification of the conditional law.}
            The preceding identities give
            \[
                X=m_Y+\xi.
            \]
            Since $Y=Tm+A\eta$ is a measurable function of $\eta$
            and $\xi\independent\eta$, we have
            \(
                \xi\independent Y
            \).
            Thus, for every bounded Borel function
            $f:E\to\mathbb R$, conditioning a measurable function
            of independent random elements gives
            \[
                \mathbb E[f(X)\mid Y]
                =
                \left.
                \int_E f(m_y+u)\,
                    \mathcal N(0,\pi_K)(du)
                \right|_{y=Y}
                \qquad\text{almost surely};
            \]
            see \cite[Lemma~8.7(i)]{kallenberg2021foundations}.
            The map $y\mapsto m_y$ is continuous on the closed affine
            subspace $Tm+T(F)$. Hence
            \[
                y\longmapsto\mathcal N(m_y,\pi_K)
            \]
            defines a Borel probability kernel on this subspace,
            which we extend by an arbitrary fixed probability measure outside
            this affine subspace.
            The displayed conditional-expectation identity identifies
            this kernel as a version of the conditional distribution
            of $X$ given $Y$. By the almost-everywhere uniqueness of
            conditional distributions
            \cite[Theorem~8.5]{kallenberg2021foundations},
            \[
                \mathscr L(X\mid T(X)=y)
                =
                \mathcal N(m_y,\pi_K)
            \]
            for $\mathscr L(T(X))$-almost every $y$.
        
            Finally, if $m\perp F$, then $m\perp K$ because
            $K\subseteq F$. Together with $m_y-m\perp K$, this gives
            $m_y\perp K$.
        \end{enumerate}
        \end{proof}

        The following proposition gives the conditional law of $U$
        at each algorithmic half-step. Its conditional mean is
        determined by the observations revealed so far, while the
        remaining Gaussian fluctuation is projected onto the
        unobserved directions.

        \begin{proposition}[Conditional Gaussian matrix laws]
        \label{proposition:mf-proof-conditional-matrix}
        For every $0\leq t\leq T$, the conditional laws of $U$ satisfy
        \begin{align}
            U&\stackrel d=_{\mathcal F_t}
            M_t+\Pi_{G_{t-1}}^\perp\check U\Pi_{W_{t-1}}^\perp,
            \label{eq:mf-proof-matrix-forward}\\
            U&\stackrel d=_{\mathcal F_t^+}
            M_t^++\Pi_{G_{t-1}}^\perp\check U\Pi_{W_t}^\perp,
            \label{eq:mf-proof-matrix-backward}
        \end{align}
        where $\check U$ is a standard Gaussian matrix independent of the
        conditioning history, and
        \begin{align}
            M_t&:= \mathbb{E}[U \mid \mathcal{F}_t] \hspace{3pt} = U\Pi_{W_{t-1}}+\Pi_{G_{t-1}}U
                -\Pi_{G_{t-1}}U\Pi_{W_{t-1}} ,
            \label{eq:mf-proof-mean-forward}\\
            M_t^+&:= \mathbb{E}[U \mid \mathcal{F}_t^+] =U\Pi_{W_t}+\Pi_{G_{t-1}}U
                -\Pi_{G_{t-1}}U\Pi_{W_t}  .
            \label{eq:mf-proof-mean-backward}
        \end{align}
        \end{proposition}

        \begin{proof}
        Equip $E:=\mathbb R^{n\times d}$ with the Frobenius inner
        product
        \(
            \langle A,B\rangle_{\mathrm F}
            :=\operatorname{tr}(A^\top B)
        \).
        We first record the geometry of the linear observations,
        then apply Lemma~\ref{lemma:mf-proof-linear-conditioning}
        successively along the algorithmic history.
        
        For matrices $W$ and $G$ with $d$ and $n$ rows, respectively,
        define
        \[
            \mathcal S(W,G)
            :=
            \{A\in E:AW=0,\ A^\top G=0\}.
        \]
        Its orthogonal projector is
        \[
            \pi_{\mathcal{S}(W,G)}(A)
            :=
            \Pi_G^\perp A\Pi_W^\perp.
        \]
        Indeed, $\pi_{\mathcal{S}(W,G)}$ is self-adjoint and idempotent
        for the Frobenius inner product. Its range is
        $\mathcal S(W,G)$: its image satisfies the two vanishing
        constraints, and every $A\in\mathcal S(W,G)$ satisfies
        \[
            A=\Pi_G^\perp A\Pi_W^\perp.
        \]
        For each realization of $U$, the affine space of matrices
        satisfying the corresponding observations is
        \[
            \mathcal A(W,G)
            :=
            \{A\in E:AW=UW,\ A^\top G=U^\top G\}
            =
            U+\mathcal S(W,G).
        \]
        Its unique element orthogonal to $\mathcal S(W,G)$ is
        \[
        \begin{aligned}
            M(W,G)
            :=
            U-\pi_{\mathcal{S}(W,G)}(U)
            =
            U\Pi_W+\Pi_GU-\Pi_GU\Pi_W.
        \end{aligned}
        \]
        To see uniqueness, any two elements of $\mathcal A(W,G)$
        differ by an element of $\mathcal S(W,G)$; if both are
        orthogonal to this subspace, their difference is zero.
        Moreover,
        \[
            M(W,G)
            =
            (UW)W^\dagger
            +(G^\top)^\dagger(U^\top G)^\top
            -\Pi_G(UW)W^\dagger.
        \]
        Thus $M(W,G)$ is a measurable function of
        $W,G,UW,U^\top G$, including when $W$ or $G$ is rank deficient.
        
        By the trajectory definitions \eqref{eq:finite-dim-macro},
        \[
            UW_s=\sqrt d\,Q_s,
            \qquad
            U^\top G_s=\sqrt n\,P_s.
        \]
        Lemma~\ref{lemma:equivalent-algorithmic-filtrations}
        therefore shows that
        \[
            M(W_{t-1},G_{t-1})
            \quad\text{and}\quad
            M(W_t,G_{t-1})
        \]
        are measurable with respect to $\mathcal F_t$ and
        $\mathcal F_t^+$, respectively. The same lemma gives
        \[
            \mathcal F_t^+
            =
            \mathcal F_t\vee\sigma(Uw^t),
            \qquad
            \mathcal F_{t+1}
            =
            \mathcal F_t^+\vee\sigma(U^\top g^t),
        \]
        where the second identity is used only for $t<T$.
        
        We use the following conditional version of
        Lemma~\ref{lemma:mf-proof-linear-conditioning}.
        Let $\mathcal G$ denote the current history, and suppose that,
        conditionally on $\mathcal G$, a random element $X$ in a
        finite-dimensional Euclidean space has law $\mathcal N(m,P)$,
        where $P$ is an orthogonal projector.
        Let $T$ be the next linear observation map.
        Assume that $m,P,T$ are $\mathcal G$-measurable, and set
        $A:=TP$.
        Then
        \[
            \mathscr L(X\mid\mathcal G\vee\sigma(TX))
            =
            \mathcal N\!\left(
                m+A^\dagger(TX-Tm),\,
                P-A^\dagger A
            \right)
            \qquad\text{almost surely}.
        \]
        Indeed, for fixed $m,P,T$, this is the Gaussian conditioning
        formula of Lemma~\ref{lemma:mf-proof-linear-conditioning}:
        $P-A^\dagger A$ is the orthogonal projector onto
        $\operatorname{ran}(P)\cap\ker T$.
        The Moore--Penrose inverse is Borel measurable as a function
        of a matrix, so
        \[
            (m,P,T,y)\longmapsto
            \mathcal N\!\left(
                m+(TP)^\dagger(y-Tm),\,
                P-(TP)^\dagger TP
            \right)
        \]
        is a measurable probability kernel.
        Integrating the deterministic conditioning identity against
        the law of the history proves the conditional-expectation
        identity for products of bounded measurable test functions.
        A monotone-class argument then identifies the conditional
        law given $\mathcal G\vee\sigma(TX)$.
        This also handles the exceptional observation sets, which
        are null under the corresponding conditional observation law. 
        The results invoked above can be found in 
        \cite[Chapter~8]{kallenberg2021foundations}.

        We now prove the conditional-law identities by induction,
        applying this formula at each half-step with the current
        history and its measurable observation map.
        
        \begin{enumerate}[wide]
            \item \emph{Initialization.}
            Since the initialization is deterministic,
            $\mathcal F_0$ is trivial.
            The matrices $W_{-1}$ and $G_{-1}$ are empty, so
            \[
                \mathcal S(W_{-1},G_{-1})=E,
                \qquad
                M(W_{-1},G_{-1})=0_E.
            \]
            Hence the forward identity at $t=0$ is precisely
            the unconditional standard Gaussian law of $U$.
        
            \item \emph{Forward conditioning.}
            Suppose the forward identity holds at time $t$.
            Given $\mathcal F_t$, the matrix $U$ is Gaussian
            with mean
            \[
                M_t=M(W_{t-1},G_{t-1})
            \]
            and covariance operator $\pi_{\mathcal S({W_{t-1},G_{t-1}})}$.
            Its conditional affine support is therefore
            \[
                M_t+\mathcal S(W_{t-1},G_{t-1})
                =
                \mathcal A(W_{t-1},G_{t-1}).
            \]
            Since $w^t$ is $\mathcal F_t$-measurable, adjoining
            $Uw^t$ amounts, conditionally on $\mathcal F_t$,
            to observing the deterministic linear map
            \[
                A\longmapsto Aw^t.
            \]
            The remaining fluctuation subspace is
            \[
            \begin{aligned}
                \mathcal S(W_{t-1},G_{t-1})
                \cap\{A:Aw^t=0\}
                =
                \mathcal S(W_t,G_{t-1}),
            \end{aligned}
            \]
            and the refined affine constraint set is
            \(
                \mathcal A(W_t,G_{t-1})
            \).
            Furthermore,
            \[
                M_t\perp\mathcal S(W_{t-1},G_{t-1}).
            \]
            The final assertion of
            Lemma~\ref{lemma:mf-proof-linear-conditioning}
            therefore shows that the new conditional mean
            is orthogonal to $\mathcal S(W_t,G_{t-1})$.
            By the uniqueness established above, this mean is
            \[
                M_t^+
                =
                M(W_t,G_{t-1}).
            \]
            The conditional covariance is the orthogonal
            projector $\pi_{\mathcal S(W_t,G_{t-1})}$.
            Since
            $\mathcal F_t^+=\mathcal F_t\vee\sigma(Uw^t)$,
            we obtain
            \[
                U\stackrel d=_{\mathcal F_t^+}
                M_t^+
                +\Pi_{G_{t-1}}^\perp
                    \check U\Pi_{W_t}^\perp.
            \]
            This proves the backward conditional-law identity
            at time $t$, including $t=T$.
        
            \item \emph{Backward conditioning.}
            Let $t<T$.
            Given $\mathcal F_t^+$, the preceding step gives
            a Gaussian law with mean $M_t^+$, fluctuation
            subspace $\mathcal S(W_t,G_{t-1})$, and affine support
            $\mathcal A(W_t,G_{t-1})$.
        
            Since $g^t$ is $\mathcal F_t^+$-measurable,
            adjoining $U^\top g^t$ amounts to observing
            the deterministic linear map
            \[
                A\longmapsto A^\top g^t.
            \]
            Its kernel intersects the current fluctuation
            subspace in
            \[
            \begin{aligned}
                \mathcal S(W_t,G_{t-1})
                \cap\{A:A^\top g^t=0\}
                =
                \mathcal S(W_t,G_t).
            \end{aligned}
            \]
            The refined affine constraint set is
            $\mathcal A(W_t,G_t)$.
            Applying
            Lemma~\ref{lemma:mf-proof-linear-conditioning}
            again, and using
            \[
                M_t^+\perp\mathcal S(W_t,G_{t-1}),
            \]
            shows that the new conditional mean is the unique
            element of $\mathcal A(W_t,G_t)$ orthogonal to
            $\mathcal S(W_t,G_t)$. It is therefore
            \[
                M_{t+1}
                =
                M(W_t,G_t).
            \]
            The new conditional covariance is
            $\pi_{\mathcal S(W_t,G_t)}$.
            Since
            $\mathcal F_{t+1}
                =\mathcal F_t^+\vee\sigma(U^\top g^t)$,
            this gives
            \[
                U\stackrel d=_{\mathcal F_{t+1}}
                M_{t+1}
                +\Pi_{G_t}^\perp\check U\Pi_{W_t}^\perp,
            \]
            which is the forward identity at time $t+1$
            and closes the induction.
        \end{enumerate}
        
        For each conditional law, the projected standard Gaussian
        matrix has the required covariance because
        $\pi_{\mathcal{S}(W,G)}$ is an orthogonal projector.
        The auxiliary matrix $\check U$ may be chosen independently
        of the conditioning history on a probability-space extension,
        separately for each representation.
        
        Finally, $\mathbb E\|U\|_{\mathrm F}^2=nd<\infty$.
        The means of the conditional Gaussian laws are consequently
        the conditional expectations of $U$, giving
        \[
            M_t=\mathbb E[U\mid\mathcal F_t],
            \qquad
            M_t^+=\mathbb E[U\mid\mathcal F_t^+].
        \]
        Expanding $M(W_{t-1},G_{t-1})$ and
        $M(W_t,G_{t-1})$ yields the stated formulas.
        \end{proof}


\subsection{Proof of the master recursion}
\label{subsec:proof-master-rec}

\begin{proof}[Proof of Proposition~\ref{proposition:master-recursion}]
Fix $T\in\mathbb N$. All limits below are taken as $d\to\infty$,
with the iteration index fixed.
Lemma~\ref{lemma:mf-proof-nondegeneracy} establishes
well-definedness, finite second moments, and absolute continuity
of the effective scores. It also gives
\[
    \bar C_t^{[w]}\succ0,
    \qquad
    \bar C_t^{[g]}\succ0,
    \qquad
    \bar v_t^{[w]}>0,
    \qquad
    \bar v_t^{[g]}>0
\]
at every finite time.
The measurability assertions in $(\mathsf F_t)$ and $(\mathsf B_t)$
follow from
Lemma~\ref{lemma:equivalent-algorithmic-filtrations}.
Their conditional matrix laws follow from
Proposition~\ref{proposition:mf-proof-conditional-matrix},
proved by successive applications of
Lemma~\ref{lemma:mf-proof-linear-conditioning}.
Thus, item~(i) of $(\mathsf F_t)$ and $(\mathsf B_t)$
already hold at every applicable time.

It remains to prove items~(ii)--(iv) of $(\mathsf F_t)$ and $(\mathsf B_t)$. We establish these assertions by alternating induction over
the forward and backward half-steps.
We freely reorder or discard columns of empirical arrays:
coordinate permutations are isometries and coordinate projections
are $1$-Lipschitz, so
Lemma~\ref{lemma:wasserstein-lipschitz-pushforward}
preserves the required Wasserstein convergence.

\begin{enumerate}[wide]
    \item \emph{Induction scheme and initialization.}
    We prove the forward and backward assertions by induction on $t$.
    Fix $0\leq t\leq T$ and assume that
    \[
        (\mathsf F_s)\ \text{and}\ (\mathsf B_s)
        \quad\text{hold for every }0\leq s<t.
    \]
    Under this induction hypothesis, we prove
    $(\mathsf F_t)$ and then, when $t<T$, $(\mathsf B_t)$.
    At $t=T$, only the forward assertion remains to be proved.
    For $t\geq1$, assertion $(\mathsf B_{t-1})$(iv) gives
    \begin{equation}
        W_2\left(
            \nu(\mu,W_t,P_{t-1}),
            \mathscr L(\tilde\mu,\tilde W_t,\tilde P_{t-1})
        \right)
        \xrightarrow{\mathbb P}0,
        \label{eq:master-proof-parameter-input}
    \end{equation}
    while $(\mathsf F_{t-1})$(iv) gives
    \begin{equation}
        W_2\left(
            \nu(Y,\Delta,Q_{t-1},G_{t-1},R_{t-1}),
            \mathscr L(
                \tilde Y,\tilde\Delta,
                \tilde Q_{t-1},\tilde G_{t-1},\tilde R_{t-1}
            )
        \right)
        \xrightarrow{\mathbb P}0
        \label{eq:master-proof-sample-input}
    \end{equation}
    and
    \begin{equation}
        b^t\xrightarrow{\mathbb P}\bar b^t.
        \label{eq:master-proof-intercept-input}
    \end{equation}
    At $t=0$, the induction hypothesis is empty.
    Since trajectory matrices with terminal index $-1$ are empty,
    \eqref{eq:master-proof-parameter-input} and
    \eqref{eq:master-proof-sample-input} reduce to the assumed
    limits of $\nu(\mu,w^0)$ and $\nu(Y,\Delta)$, respectively.
    Equation~\eqref{eq:master-proof-intercept-input} follows
    from the initialization assumption.
    Thus these three inputs are available at every induction
    step, including $t=0$.
    
    \item \emph{Forward half-step $(\mathsf F_t)$.} Let $0\leq t\leq T$.
    Under the induction hypothesis above, we establish
    items~(ii)--(iv) of $(\mathsf F_t)$ using
    \eqref{eq:master-proof-parameter-input}--%
    \eqref{eq:master-proof-intercept-input}.
    
    \begin{enumerate}[label=(\alph*),wide]
        \item \emph{Projection coefficients and nondegeneracy $(\mathsf F_t)$\textup{(ii)}.}
        Part~1 of
        Lemma~\ref{lemma:mf-proof-empirical-stability}
        applied to the parameter-side input gives
        \[
            C_t^{[w]}
            \xrightarrow{\mathbb P}\bar C_t^{[w]},
            \qquad
            m^t
            =
            \frac{\mu^\top w^t}{d}
            \xrightarrow{\mathbb P}
            \tilde{\mathbb E}[\tilde\mu\tilde w^t]
            =
            \bar m^t.
        \]
        For $t\geq1$, it also gives
        \[
            \frac{W_{t-1}^\top w^t}{d}
            \xrightarrow{\mathbb P}
            \left(
                \tilde{\mathbb E}[\tilde w^s\tilde w^t]
            \right)_{0\leq s<t}.
        \]
        Since $\bar C_{t-1}^{[w]}\succ0$, the
        pseudoinverse assertion of the same lemma yields
        \(
            \alpha_t^{[w]}
            \xrightarrow{\mathbb P}\bar\alpha_t^{[w]}
        \).
        The limiting vector solves the weight projection
        equations in
        Definition~\ref{def:mean-field-recursion}.
        Apply part~2 of
        Lemma~\ref{lemma:mf-proof-empirical-stability}
        to the parameter-side array
        \[
            H^{(d)}
            :=
            [\,\mu\mid W_t\mid P_{t-1}\,]
            \in\mathbb R^{d\times(2t+2)},
        \]
        whose empirical convergence is
        \eqref{eq:master-proof-parameter-input},
        and the coefficient vector $\alpha_t^{[w]}$.
        The rowwise map
        \[
            F\bigl(\alpha,(u,w_0,\ldots,w_t,p_0,\ldots,p_{t-1})\bigr)
            :=
            w_t-\sum_{s=0}^{t-1}\alpha_s w_s
        \]
        is continuous and has linear growth uniformly for $\alpha$
        near $\bar\alpha_t^{[w]}$.
        Its application to $H^{(d)}$ gives
        \(
            w_\perp^t=w^t-W_{t-1}\alpha_t^{[w]}
        \).
        Hence Lemma~\ref{lemma:mf-proof-empirical-stability} yields
        \[
            W_2\left(
                \nu(\mu,W_t,P_{t-1},w_\perp^t),
                \mathscr L(
                    \tilde\mu,\tilde W_t,\tilde P_{t-1},
                    \tilde w_\perp^t
                )
            \right)
            \xrightarrow{\mathbb P}0.
        \]
        Taking second moments and cross moments then gives
        \[
        \begin{aligned}
            \frac{\|w_\perp^t\|^2}{d}
            \xrightarrow{\mathbb P}
            \tilde{\mathbb E}[(\tilde w_\perp^t)^2]
            , \qquad
            \frac{P_{t-1}^\top w_\perp^t}{d}
            \xrightarrow{\mathbb P}
            \left(
                \tilde{\mathbb E}[\tilde p^s\tilde w_\perp^t]
            \right)_{0\leq s<t}.
        \end{aligned}
        \]
        The sample-side input similarly implies
        \[
            C_{t-1}^{[g]}
            \xrightarrow{\mathbb P}\bar C_{t-1}^{[g]}.
        \]
        Its limit is positive definite when $t\geq1$.
        Consequently,
        \[
            \phi_t^{[w]}
            \xrightarrow{\mathbb P}\bar\phi_t^{[w]},
            \qquad
            v_t^{[w]}
            \xrightarrow{\mathbb P}\bar v_t^{[w]}>0.
        \]
        The first limit satisfies the defining forward
        memory equations, and the second follows by taking
        nonnegative square roots.
        At $t=0$, both coefficient vectors are empty and
        \[
            (v_0^{[w]})^2
            =
            \frac{\|w^0\|^2}{d}
            \longrightarrow
            \tilde{\mathbb E}[(\tilde w^0)^2]
            =
            (\bar v_0^{[w]})^2>0.
        \]
        Thus these conclusions also hold at initialization.
        This proves $(\mathsf F_t)$(ii), as well as the convergence of $m^t$.

        \item \emph{Conditional forward decomposition and error $(\mathsf F_t)$\textup{(iii)}.}
        The projection definitions give
        \[
            \Pi_{W_{t-1}}w^t
            =
            W_{t-1}\alpha_t^{[w]},
            \qquad
            w_\perp^t
            =
            \Pi_{W_{t-1}}^\perp w^t.
        \]
        Since $w^t$ is $\mathcal F_t$-measurable,
        multiplying the conditional matrix-law identity
        by $w^t/\sqrt d$ yields
        \[
            q^t
            \stackrel d=_{\mathcal F_t}
            \frac{M_tw^t}{\sqrt d}
            +
            \Pi_{G_{t-1}}^\perp
            \frac{\check Uw_\perp^t}{\sqrt d}.
        \]
        The conditional mean decomposes as
        \[
            \frac{M_tw^t}{\sqrt d}
            =
            \frac{U\Pi_{W_{t-1}}w^t}{\sqrt d}
            +
            \frac{\Pi_{G_{t-1}}Uw_\perp^t}{\sqrt d}.
        \]
        The first term equals
        $Q_{t-1}\alpha_t^{[w]}$.
        For the second, the identities
        \[
            \Pi_A=A(A^\top A)^\dagger A^\top,
            \qquad
            G_{t-1}^\top U=\sqrt n\,P_{t-1}^\top
        \]
        give
        \[
        \begin{aligned}
            \frac{\Pi_{G_{t-1}}Uw_\perp^t}{\sqrt d}
            =
            \frac{\sqrt n}{\sqrt d}\,
            G_{t-1}(G_{t-1}^\top G_{t-1})^\dagger
            P_{t-1}^\top w_\perp^t
            =
            \frac1{\sqrt\delta}\,
            G_{t-1}(C_{t-1}^{[g]})^\dagger
            \frac{P_{t-1}^\top w_\perp^t}{d}
            =
            \frac1{\sqrt\delta}\,
            G_{t-1}\phi_t^{[w]}.
        \end{aligned}
        \]
        Conditionally on $\mathcal F_t$,
        $\check Uw_\perp^t/\sqrt d$ is centered Gaussian
        with covariance $(v_t^{[w]})^2I_n$.
        Identifying Gaussian laws by their means and
        covariances
        \cite[Lemma~14.1(i)]{kallenberg2021foundations},
        the fluctuation therefore has the same conditional
        law as
        \[
            v_t^{[w]}\Pi_{G_{t-1}}^\perp z_t^{[g]}
            =
            v_t^{[w]}z_t^{[g]}+\varepsilon_t^{[g]},
        \]
        where $z_t^{[g]}\sim\mathcal N(0,I_n)$ is independent
        of $\mathcal F_t$ and
        \[
            \varepsilon_t^{[g]}
            =
            -v_t^{[w]}\Pi_{G_{t-1}}z_t^{[g]}.
        \]
        This proves
        \eqref{eq:conditional-q-decomposition-lemmas},
        including when $v_t^{[w]}=0$.
        Since the projector and innovation amplitude are
        $\mathcal F_t$-measurable,
        \[
        \begin{aligned}
            \mathbb E\left[
                \frac{\|\varepsilon_t^{[g]}\|^2}{n}
                \Bigm|\mathcal F_t
            \right]
            =
            \frac{(v_t^{[w]})^2}{n}
            \operatorname{tr}(\Pi_{G_{t-1}})
            =
            \frac{(v_t^{[w]})^2}{n}
            \operatorname{rank}(G_{t-1})
            \leq
            (v_t^{[w]})^2\frac tn.
        \end{aligned}
        \]
        To obtain convergence of the error, the conditional
        Markov inequality
        gives, for $K,\epsilon>0$,
        \[
            \mathbb P\left(
                \frac{\|\varepsilon_t^{[g]}\|^2}{n}
                >\epsilon
            \right)
            \leq
            \mathbb P\bigl((v_t^{[w]})^2>K\bigr)
            +
            \frac{Kt}{n\epsilon}.
        \]
        Choosing $K>(\bar v_t^{[w]})^2$, both terms tend
        to zero. This proves $(\mathsf F_t)$(iii).

        \item \emph{Sample-side empirical extension $(\mathsf F_t)$\textup{(iv)}.}
        Apply
        Lemma~\ref{lemma:mf-proof-gaussian-extension}
        to the conditional decomposition just established,
        using the history array
        \[
            [\,Y\mid\Delta\mid Q_{t-1}
                \mid G_{t-1}\mid R_{t-1}\,].
        \]
        This array and all the coefficients are
        $\mathcal F_t$-measurable.
        Its empirical law converges by
        \eqref{eq:master-proof-sample-input},
        the coefficients converge by part~(a) and
        $\delta\to\bar\delta$, and the normalized squared
        error tends to zero by part~(b).

        The limiting new coordinate is
        \[
            \sum_{s=0}^{t-1}
                \bar\alpha^{[w]}(t,s)\tilde q^s
            +
            \frac1{\sqrt{\bar\delta}}
            \sum_{s=0}^{t-1}
                \bar\phi^{[w]}(t,s)\tilde g^s
            +
            \bar v_t^{[w]}z_t^{[g]}
            =
            \tilde q^t.
        \]
        Here the fresh scalar Gaussian is independent of
        the limiting history by construction.
        The equality with $\tilde q^t$ is
        \eqref{eq:mf-proof-forward-one-step}, established in the
        proof of Lemma~\ref{lemma:mf-proof-nondegeneracy}.
        Thus
        \begin{equation}
            W_2\left(
                \nu(Y,\Delta,Q_t,G_{t-1},R_{t-1}),
                \mathscr L(
                    \tilde Y,\tilde\Delta,
                    \tilde Q_t,\tilde G_{t-1},\tilde R_{t-1}
                )
            \right)
            \xrightarrow{\mathbb P}0.
            \label{eq:master-proof-forward-q-limit}
        \end{equation}
        The score relation is
        \[
            r_i^t=b^t+m^tY_i+\sigma q_i^t.
        \]
        As before, its rowwise map is continuous and has linear
        growth uniformly for $(b^t,m^t)$ near
        $(\bar b^t,\bar m^t)$.
        Applying part~2 of
        Lemma~\ref{lemma:mf-proof-empirical-stability}
        therefore appends
        \[
            \tilde r^t
            =
            \bar b^t+\bar m^t\tilde Y+\sigma\tilde q^t
        \]
        to the joint limit.
        After reordering columns, this is precisely
        \eqref{eq:master-forward-empirical-law}.
        If $t=T$, this completes $(\mathsf F_T)$.
        Suppose henceforth that $t<T$.
        Define, for $u\in(0,1)$, $y\in\{-1,+1\}$,
        $a\in\{0,1\}$, and $r\in\mathbb R$,
        \[
            \mathcal R(u,y,a,r)
            :=\frac a u\,\ell'(y,r)
            +\frac{\pi(1-a)}{1-u}
                S_\kappa(r)\ell'(\operatorname{sign}(r),r).
        \]
        Then \eqref{eq:residual-map} gives
        $g_i^t=\mathcal R(\rho,Y_i,\Delta_i,r_i^t)$.
        On the finite label and supervision state space,
        this map is jointly continuous at every point
        with $\rho\in(0,1)$ and
        $r\notin\{\kappa_-,\kappa_+\}$.
        The inactive interval contains zero, so the
        pseudo-label convention at zero introduces no
        additional discontinuity in the residual.

        The Lipschitz loss derivatives and bounded
        selection indicators give, for some $C<\infty$,
        \[
            |\mathcal R(\varrho,y,a,r)|
            \leq C(1+|r|)
        \]
        uniformly over $\varrho$ in a neighborhood of
        $\bar\rho$, $y\in\{-1,+1\}$,
        $a\in\{0,1\}$, and $r\in\mathbb R$.        
        We apply the stability Lemma~\ref{lemma:mf-proof-empirical-stability} on the closed domain
        \[
            \mathcal X_t
            :=
            \{-1,+1\}\times\{0,1\}
            \times\mathbb R^{t+1}
            \times\mathbb R^t
            \times\mathbb R^{t+1},
        \]
        corresponding to the columns of
        \[
            [\,Y\mid\Delta\mid Q_t\mid G_{t-1}\mid R_t\,].
        \]
        The empirical rows and the limiting history take values
        in $\mathcal X_t$.
        In its relative topology, the residual map is continuous
        at every point with $\varrho\in(0,1)$ and current score
        outside $\{\kappa_-,\kappa_+\}$.
        By Lemma~\ref{lemma:mf-proof-nondegeneracy},
        \[
            \tilde{\mathbb P}
            (\tilde r^t\in\{\kappa_-,\kappa_+\})=0.
        \]
        Together with $\rho\to\bar\rho\in(0,1)$ and the uniform
        linear-growth bound, this verifies the hypotheses of
        part~2 of Lemma~\ref{lemma:mf-proof-empirical-stability}.
        Applying that lemma appends $g^t$ and proves
        \eqref{eq:master-forward-residual-extension}.
        Part~1 of the same lemma now gives
        \[
            \zeta^t
            =
            \frac1n1_n^\top g^t
            \xrightarrow{\mathbb P}
            \tilde{\mathbb E}[\tilde g^t]
            =
            \bar\zeta^t,
            \qquad
            \chi^t
            =
            \frac1nY^\top g^t
            \xrightarrow{\mathbb P}
            \tilde{\mathbb E}[\tilde Y\tilde g^t]
            =
            \bar\chi^t.
        \]
        Together with the intercept input, the update
        \(
            b^{t+1}=b^t-\eta\zeta^t
        \)
        yields
        $b^{t+1}\xrightarrow{\mathbb P}\bar b^{t+1}$.
        This proves $(\mathsf F_t)$(iv), and hence
        $(\mathsf F_t)$ which concludes the forward half-step.
    \end{enumerate}

    \item \emph{Backward half-step $(\mathsf B_t)$.}
    Let $0\leq t<T$, and suppose $(\mathsf F_t)$ has been
    established. We retain the parameter-side induction input
    \eqref{eq:master-proof-parameter-input}.
    \begin{enumerate}[label=(\alph*),wide]
        \item \emph{Projection coefficients and nondegeneracy
        $(\mathsf B_t)$\textup{(ii)}.}
        The residual extension
        \eqref{eq:master-forward-residual-extension}
        and part~1 of
        Lemma~\ref{lemma:mf-proof-empirical-stability}
        give
        \[
            C_t^{[g]}
            \xrightarrow{\mathbb P}\bar C_t^{[g]}\succ0.
        \]
        For $t\geq1$, they also give
        \[
            \frac{G_{t-1}^\top g^t}{n}
            \xrightarrow{\mathbb P}
            \left(
                \tilde{\mathbb E}[\tilde g^s\tilde g^t]
            \right)_{0\leq s<t}.
        \]
        The leading principal block $C_{t-1}^{[g]}$
        converges to $\bar C_{t-1}^{[g]}\succ0$.
        The pseudoinverse assertion of the same lemma
        therefore yields
        \(
            \alpha_t^{[g]}
            \xrightarrow{\mathbb P}\bar\alpha_t^{[g]}
        \).
        The limiting vector solves the residual projection
        equations in
        Definition~\ref{def:mean-field-recursion}.
        At $t=0$, this coefficient vector is empty.
        Apply part~2 of
        Lemma~\ref{lemma:mf-proof-empirical-stability}
        to the sample-side array
        \[
            H^{(n)}
            :=
            [\,Y\mid\Delta\mid Q_t\mid G_t\mid R_t\,]
            \in\mathbb R^{n\times(3t+5)},
        \]
        whose empirical convergence is
        \eqref{eq:master-forward-residual-extension},
        and the coefficient vector $\alpha_t^{[g]}$.
        The rowwise map
        \[
        \begin{aligned}
            F\bigl(
                \alpha,
                (y,a,q_0,\ldots,q_t,
                    g_0,\ldots,g_t,r_0,\ldots,r_t)
            \bigr)
            :=
            g_t-\sum_{s=0}^{t-1}\alpha_s g_s
        \end{aligned}
        \]
        is continuous and has linear growth uniformly for
        $\alpha$ near $\bar\alpha_t^{[g]}$.
        Its application to $H^{(n)}$ gives
        \[
            g_\perp^t=g^t-G_{t-1}\alpha_t^{[g]}.
        \]
        Hence the stability lemma yields
        \[
            W_2\left(
                \nu(Y,\Delta,Q_t,G_t,R_t,g_\perp^t),
                \mathscr L(
                    \tilde Y,\tilde\Delta,\tilde Q_t,
                    \tilde G_t,\tilde R_t,\tilde g_\perp^t
                )
            \right)
            \xrightarrow{\mathbb P}0.
        \]
        Taking second moments and cross moments then gives
        \[
        \begin{aligned}
            \frac{\|g_\perp^t\|^2}{n}
            \xrightarrow{\mathbb P}
            \tilde{\mathbb E}[(\tilde g_\perp^t)^2],\qquad
            \frac{Q_t^\top g_\perp^t}{n}
            \xrightarrow{\mathbb P}
            \left(
                \tilde{\mathbb E}[\tilde q^s\tilde g_\perp^t]
            \right)_{0\leq s\leq t}.
        \end{aligned}
        \]
        By \eqref{eq:master-forward-nondegeneracy},
        \(
            C_t^{[w]}
            \xrightarrow{\mathbb P}\bar C_t^{[w]}\succ0
        \).
        Consequently,
        \[
            \phi_t^{[g]}
            \xrightarrow{\mathbb P}\bar\phi_t^{[g]},
            \qquad
            v_t^{[g]}
            \xrightarrow{\mathbb P}\bar v_t^{[g]}>0.
        \]
        The first limit satisfies the defining backward
        memory equations, and the second follows by taking
        nonnegative square roots.
        At $t=0$, the same argument applies with $g_\perp^0=g^0$
        and the empty coefficient vector $\alpha_0^{[g]}$;
        the memory coefficient vector $\phi_0^{[g]}$ need not be empty.
        This proves $(\mathsf B_t)$(ii).
    
        \item \emph{Conditional backward decomposition and error
        $(\mathsf B_t)$\textup{(iii)}.}
        The projection definitions give
        \[
            \Pi_{G_{t-1}}g^t
            =
            G_{t-1}\alpha_t^{[g]},
            \qquad
            g_\perp^t
            =
            \Pi_{G_{t-1}}^\perp g^t.
        \]
        Since $g^t$ is $\mathcal F_t^+$-measurable,
        transposing \eqref{eq:mf-proof-matrix-backward}
        and multiplying by $g^t/\sqrt n$ yields
        \[
            p^t
            \stackrel d=_{\mathcal F_t^+}
            \frac{(M_t^+)^\top g^t}{\sqrt n}
            +
            \Pi_{W_t}^\perp
            \frac{\check U^\top g_\perp^t}{\sqrt n}.
        \]
        Using \eqref{eq:mf-proof-mean-backward},
        the conditional mean decomposes as
        \[
            \frac{(M_t^+)^\top g^t}{\sqrt n}
            =
            \frac{U^\top\Pi_{G_{t-1}}g^t}{\sqrt n}
            +
            \frac{\Pi_{W_t}U^\top g_\perp^t}{\sqrt n}.
        \]
        Since
        $U^\top G_{t-1}=\sqrt n\,P_{t-1}$,
        the first term equals
        $P_{t-1}\alpha_t^{[g]}$.
        For the second, the identities
        \[
            \Pi_A=A(A^\top A)^\dagger A^\top,
            \qquad
            W_t^\top U^\top=\sqrt d\,Q_t^\top
        \]
        give
        \[
        \begin{aligned}
            \frac{\Pi_{W_t}U^\top g_\perp^t}{\sqrt n}
            =
            \frac{\sqrt d}{\sqrt n}\,
            W_t(W_t^\top W_t)^\dagger
            Q_t^\top g_\perp^t
            =
            \sqrt\delta\,
            W_t(C_t^{[w]})^\dagger
            \frac{Q_t^\top g_\perp^t}{n}
            =
            \sqrt\delta\,W_t\phi_t^{[g]}.
        \end{aligned}
        \]
        Conditionally on $\mathcal F_t^+$,
        $\check U^\top g_\perp^t/\sqrt n$ is centered
        Gaussian with covariance
        \[
            \frac{\|g_\perp^t\|^2}{n}I_d
            =
            (v_t^{[g]})^2I_d.
        \]
        Identifying Gaussian laws by their means and
        covariances
        \cite[Lemma~14.1(i)]{kallenberg2021foundations},
        its projected version has the same conditional law as
        \[
            v_t^{[g]}\Pi_{W_t}^\perp z_t^{[w]}
            =
            v_t^{[g]}z_t^{[w]}+\varepsilon_t^{[w]},
        \]
        where $z_t^{[w]}\sim\mathcal N(0,I_d)$ is independent
        of $\mathcal F_t^+$ and
        \[
            \varepsilon_t^{[w]}
            =
            -v_t^{[g]}\Pi_{W_t}z_t^{[w]}.
        \]
        This proves
        \eqref{eq:conditional-p-decomposition-lemmas},
        including when $v_t^{[g]}=0$.
        Since the projector and innovation amplitude are
        $\mathcal F_t^+$-measurable,
        \[
        \begin{aligned}
            \mathbb E\left[
                \frac{\|\varepsilon_t^{[w]}\|^2}{d}
                \Bigm|\mathcal F_t^+
            \right]
            =
            \frac{(v_t^{[g]})^2}{d}
            \operatorname{tr}(\Pi_{W_t})
            =
            \frac{(v_t^{[g]})^2}{d}
            \operatorname{rank}(W_t)
            \leq
            (v_t^{[g]})^2\frac{t+1}{d}.
        \end{aligned}
        \]
        Thus \eqref{eq:backward-projection-error-lemmas} holds. The
        conditional Markov inequality gives, for
        $K,\epsilon>0$,
        \[
            \mathbb P\left(
                \frac{\|\varepsilon_t^{[w]}\|^2}{d}
                >\epsilon
            \right)
            \leq
            \mathbb P\bigl((v_t^{[g]})^2>K\bigr)
            +
            \frac{K(t+1)}{d\epsilon}.
        \]
        Choosing $K>(\bar v_t^{[g]})^2$, both terms tend
        to zero. This proves $(\mathsf B_t)$(iii).
    
        \item \emph{Parameter-side empirical extension
        $(\mathsf B_t)$\textup{(iv)}.}
        Apply
        Lemma~\ref{lemma:mf-proof-gaussian-extension}
        to the conditional decomposition just established,
        using the parameter-side history array
        \[
            H^{(d)}
            :=
            [\,\mu\mid W_t\mid P_{t-1}\,]
            \in\mathbb R^{d\times(2t+2)}.
        \]
        This array and all the coefficients are
        $\mathcal F_t^+$-measurable.
        Its empirical law converges by
        \eqref{eq:master-proof-parameter-input},
        the coefficients converge by part~(a) and
        $\delta\to\bar\delta$, and the normalized squared
        error tends to zero by part~(b).
        The limiting new coordinate is
        \[
            \sum_{s=0}^{t-1}
                \bar\alpha^{[g]}(t,s)\tilde p^s
            +
            \sqrt{\bar\delta}
            \sum_{s=0}^{t}
                \bar\phi^{[g]}(t,s)\tilde w^s
            +
            \bar v_t^{[g]}z_t^{[w]}
            =
            \tilde p^t.
        \]
        Here the fresh scalar Gaussian is independent of
        the limiting history by construction.
        The equality with $\tilde p^t$ is
        \eqref{eq:mf-proof-backward-one-step}, established in the
        proof of Lemma~\ref{lemma:mf-proof-nondegeneracy}.
        Hence
        \begin{equation}
            W_2\left(
                \nu(\mu,W_t,P_t),
                \mathscr L(\tilde\mu,\tilde W_t,\tilde P_t)
            \right)
            \xrightarrow{\mathbb P}0.
            \label{eq:master-proof-backward-p-limit}
        \end{equation}
        We next append the weight update.
        Apply part~2 of
        Lemma~\ref{lemma:mf-proof-empirical-stability}
        to the augmented parameter-side array
        \[
            H^{(d)}_+
            :=
            [\,\mu\mid W_t\mid P_t\,]
            \in\mathbb R^{d\times(2t+3)}
        \]
        and the coefficient vector
        \[
            \beta^{(d)}
            :=
            \left(\chi^t,\frac{\sigma}{\sqrt\delta}\right)
            \xrightarrow{\mathbb P}
            \left(\bar\chi^t,\frac{\sigma}{\sqrt{\bar\delta}}\right).
        \]
        This coefficient convergence follows from
        \eqref{eq:master-backward-macroscopic-convergence}
        and $\delta\to\bar\delta>0$.
        The rowwise map
        \[
        \begin{aligned}
            F\bigl(
                (a,c),
                (u,w_0,\ldots,w_t,p_0,\ldots,p_t)
            \bigr)
            :=
            w_t-\eta\left[
                \lambda j'(w_t)+au+cp_t
            \right]
        \end{aligned}
        \]
        is continuous and has linear growth uniformly for
        $(a,c)$ near
        $(\bar\chi^t,\sigma/\sqrt{\bar\delta})$,
        because $j'$ is Lipschitz.
        Its application to $H^{(d)}_+$ gives
        \[
            w_k^{t+1}
            =
            w_k^t-\eta\left[
                \lambda j'(w_k^t)
                +\chi^t\mu_k
                +\frac{\sigma}{\sqrt\delta}p_k^t
            \right].
        \]
        Thus the stability Lemma~\ref{lemma:mf-proof-empirical-stability} appends the effective coordinate
        \[
            \tilde w^{t+1}
            =
            \tilde w^t-\eta\left[
                \lambda j'(\tilde w^t)
                +\bar\chi^t\tilde\mu
                +\frac{\sigma}{\sqrt{\bar\delta}}\tilde p^t
            \right]
        \]
        to the joint limit
        \eqref{eq:master-proof-backward-p-limit}.
        After reordering columns, this proves
        \eqref{eq:master-backward-empirical-law}.
        This establishes $(\mathsf B_t)$(iv), and hence
        $(\mathsf B_t)$.
    \end{enumerate}
\end{enumerate}

We have shown that the assertions at times $s<t$ imply
$(\mathsf F_t)$ and, when $t<T$, $(\mathsf B_t)$.
The same arguments apply at $t=0$ using the initialization
limits established above.
Strong induction therefore proves $(\mathsf F_t)$ for
$0\leq t\leq T$ and $(\mathsf B_t)$ for $0\leq t<T$.

\end{proof}

\subsection{Proof of Theorem~\ref{theorem:trajectory}
and Corollary~\ref{cor:mean-field-observables}}
\label{subsec:proof-trajectory}

\begin{proof}[Proof of Theorem~\ref{theorem:trajectory}]
Well-definedness, uniqueness, finite second moments, and
atomlessness of the effective scores follow from
Lemma~\ref{lemma:mf-proof-nondegeneracy}.
Fix a finite horizon $T\in\mathbb N$.
The remaining assertions follow by projecting the joint
empirical limits in
Proposition~\ref{proposition:master-recursion}.

By Lemma~\ref{lemma:wasserstein-lipschitz-pushforward},
coordinate projections preserve $W_2$ convergence in probability.
\begin{enumerate}[label=(\roman*),wide]
    \item \emph{Parameter-side trajectories.}
    For $T\geq1$, apply
    \eqref{eq:master-backward-empirical-law}
    at $t=T-1$ and project onto the coordinates
    $(\mu,W_T)$. This gives
    \[
        W_2\left(
            \nu(\mu,W_T),
            \mathscr L(\tilde\mu,\tilde W_T)
        \right)
        \xrightarrow{\mathbb P}0,
    \]
    which is \eqref{eq:parameter-trajectory-limit}.
    For $T=0$, the same conclusion follows directly
    from Assumption~\ref{assumption:quenched-initialization}.

    \item \emph{Sample-side trajectories.}
    Apply \eqref{eq:master-forward-empirical-law}
    at $t=T$ and project onto $(Y,\Delta,R_T)$.
    This yields
    \[
        W_2\left(
            \nu(Y,\Delta,R_T),
            \mathscr L(\tilde Y,\tilde\Delta,\tilde R_T)
        \right)
        \xrightarrow{\mathbb P}0,
    \]
    proving \eqref{eq:sample-trajectory-limit}.

    \item \emph{Intercepts.}
    The initialization assumption gives
    $b^0\to\bar b^0$.
    For every $1\leq t\leq T$,
    \eqref{eq:master-backward-macroscopic-convergence}
    applied at time $t-1$ gives
    $b^t\xrightarrow{\mathbb P}\bar b^t$.
    Hence, for every $\epsilon>0$,
    \[
        \mathbb P\left(
            \max_{0\leq t\leq T}|b^t-\bar b^t|>\epsilon
        \right)
        \leq
        \sum_{t=0}^T
        \mathbb P(|b^t-\bar b^t|>\epsilon)
        \longrightarrow0.
    \]
    This proves \eqref{eq:mean-field-bias-convergence}.
\end{enumerate}
\end{proof}

\begin{proof}[Proof of Corollary~\ref{cor:mean-field-observables}]
Fix nonnegative integers $s,t$ and choose a finite horizon
$T\geq\max\{s,t\}$.

\begin{enumerate}[wide]
    \item \emph{Weight-side observables.}
    Apply part~1 of
    Lemma~\ref{lemma:mf-proof-empirical-stability}
    to the parameter-side trajectory convergence
    \eqref{eq:parameter-trajectory-limit}.
    It gives
    \[
    \begin{aligned}
        m^t
        &=
        \frac{\mu^\top w^t}{d}
        \xrightarrow{\mathbb P}
        \tilde{\mathbb E}[\tilde\mu\tilde w^t]
        =
        \bar m^t,\\
        (\tau^t)^2
        &=
        \frac{\|w^t\|^2}{d}
        \xrightarrow{\mathbb P}
        \tilde{\mathbb E}[(\tilde w^t)^2]
        =
        (\bar\tau^t)^2,\\
        C^{[w]}(s,t)
        &=
        \frac{(w^s)^\top w^t}{d}
        \xrightarrow{\mathbb P}
        \tilde{\mathbb E}[\tilde w^s\tilde w^t]
        =
        \bar C^{[w]}(s,t).
    \end{aligned}
    \]
    Taking nonnegative square roots yields
    $\tau^t\xrightarrow{\mathbb P}\bar\tau^t$.
    Moreover,
    \[
        (\bar\tau^t)^2=\bar C^{[w]}(t,t)>0
    \]
    by Lemma~\ref{lemma:mf-proof-nondegeneracy}.

    \item \emph{Selected fraction.}
    Projecting \eqref{eq:sample-trajectory-limit}
    onto the supervision indicator and the score at
    time $t$ gives
    \[
        W_2\left(
            \nu(\Delta,r^t),
            \mathscr L(\tilde\Delta,\tilde r^t)
        \right)
        \xrightarrow{\mathbb P}0.
    \]
    Consider the empirical array
    \[
        H^{(n)}:=[\,\Delta\mid r^t\,]
        \in\mathbb R^{n\times2},
    \]
    whose rows and limiting law take values in the
    closed set $\mathcal X:=\{0,1\}\times\mathbb R$.
    Since $M=n(1-\rho)$,
    \[
        \omega^t
        =
        \frac1n\sum_{i=1}^n
        \frac{1-\Delta_i}{1-\rho}S_\kappa(r_i^t).
    \]
    Apply part~2 of
    Lemma~\ref{lemma:mf-proof-empirical-stability}
    with coefficient and rowwise map
    \[
        \frac1{1-\rho}
        \longrightarrow
        \frac1{1-\bar\rho} , \quad \text{and} \quad  F(c,(a,r)):=c(1-a)S_\kappa(r),
    \]
    respectively.
    This map is bounded uniformly for $c$ near
    $(1-\bar\rho)^{-1}$, and its only possible
    discontinuities on $\mathcal X$ occur at
    $r=\kappa_-$ and $r=\kappa_+$.
    The limiting score law is atomless by
    Lemma~\ref{lemma:mf-proof-nondegeneracy}, so
    \[
        \tilde{\mathbb P}
        (\tilde r^t\in\{\kappa_-,\kappa_+\})=0.
    \]
    Thus the stability Lemma~\ref{lemma:mf-proof-empirical-stability} applies.
    Taking the empirical first moment of the appended
    coordinate using part~1 of that lemma yields
    \[
        \omega^t
        \xrightarrow{\mathbb P}
        \frac{
            \tilde{\mathbb E}
            [(1-\tilde\Delta)S_\kappa(\tilde r^t)]
        }{1-\bar\rho}
        =
        \bar\omega^t.
    \]
    Together with part~1, this proves
    \eqref{eq:mean-field-observable-convergence}..

    \item \emph{Selected-error fraction}
    Apply the same argument
    to the array $(Y,\Delta,r^t)$ and the rowwise map
    \[
        F(c,(y,a,r))
        :=
        c(1-a)
        \mathbf1_{\{
            \mathbf1_{\{r\geq\kappa_+\}}
            -\mathbf1_{\{r\leq\kappa_-\}}
            =-y
        \}}.
    \]
    On the discrete label and supervision state space, its only
    possible discontinuities occur at the selection thresholds.
    Lemma~\ref{lemma:mf-proof-nondegeneracy} and part~2 of
    Lemma~\ref{lemma:mf-proof-empirical-stability}
    therefore give $e^t\xrightarrow{\mathbb P}\bar e^t$.

    \item \emph{Classification error.}
    Let $\mathcal F_{\mathrm{train}}$ denote the
    sigma-algebra generated by the training data.
    Draw an independent test observation
    \[
        X_{\mathrm{test}}
        =
        \frac{Y_{\mathrm{test}}\mu}{\sqrt d}
        +\sigma Z_{\mathrm{test}},
    \]
    where
    \[
        \mathbb P(Y_{\mathrm{test}}=+1)=p,
        \qquad
        Z_{\mathrm{test}}\sim\mathcal N(0,I_d),
    \]
    and the label and Gaussian noise are independent
    of each other and of $\mathcal F_{\mathrm{train}}$.
    The test score satisfies
    \[
    \begin{aligned}
        r_{\mathrm{test}}^t
        :=
        b^t+\frac{(w^t)^\top X_{\mathrm{test}}}{\sqrt d}
        =
        b^t+m^tY_{\mathrm{test}}
        +\frac{\sigma}{\sqrt d}
            (w^t)^\top Z_{\mathrm{test}}.
    \end{aligned}
    \]
    Since $b^t$ and $w^t$ are
    $\mathcal F_{\mathrm{train}}$-measurable,
    \[
        \mathscr L\left(
            r_{\mathrm{test}}^t
            \mid\mathcal F_{\mathrm{train}},
                Y_{\mathrm{test}}=y
        \right)
        =
        \mathcal N\left(
            b^t+ym^t,\,
            \sigma^2(\tau^t)^2
        \right).
    \]
    On $\{\tau^t>0\}$, this distribution assigns zero
    probability to a zero score.
    The Gaussian tail formula and symmetry of $\Phi$
    therefore give
    \[
        \mathcal E^t
        =
        p\,\Phi\left(
            \frac{-b^t-m^t}{\sigma\tau^t}
        \right)
        +
        (1-p)\,\Phi\left(
            \frac{b^t-m^t}{\sigma\tau^t}
        \right).
    \]
    Assumption~\ref{assumption:asymptotic-regime},
    Theorem~\ref{theorem:trajectory}, and part~1 give
    \[
        (p,b^t,m^t,\tau^t)
        \xrightarrow{\mathbb P}
        (\bar p,\bar b^t,\bar m^t,\bar\tau^t),
        \qquad
        \bar\tau^t>0.
    \]
    Define $\widehat{\mathcal E}^t$ by the displayed
    Gaussian expression when $\tau^t>0$, and set it
    equal to zero when $\tau^t=0$.
    This defines a measurable function of
    $(p,b^t,m^t,\tau^t)$ that is continuous at the
    limiting values.
    The continuous mapping theorem
    \cite[Lemmas~5.3--5.4]{kallenberg2021foundations}
    consequently yields
    \[
        \widehat{\mathcal E}^t
        \xrightarrow{\mathbb P}
        \bar p\,\Phi\left(
            \frac{-\bar b^t-\bar m^t}{\sigma\bar\tau^t}
        \right)
        +
        (1-\bar p)\,\Phi\left(
            \frac{\bar b^t-\bar m^t}{\sigma\bar\tau^t}
        \right)
        =
        \bar{\mathcal E}^t.
    \]
    Finally,
    \[
        \mathbb P(\tau^t=0)
        \leq
        \mathbb P\left(
            |\tau^t-\bar\tau^t|>\frac{\bar\tau^t}{2}
        \right)
        \longrightarrow0.
    \]
    Since $\mathcal E^t=\widehat{\mathcal E}^t$
    on $\{\tau^t>0\}$, for every $\epsilon>0$,
    \[
        \mathbb P\left(
            |\mathcal E^t-\bar{\mathcal E}^t|>\epsilon
        \right)
        \leq
        \mathbb P(\tau^t=0)
        +
        \mathbb P\left(
            |\widehat{\mathcal E}^t-\bar{\mathcal E}^t|
            >\epsilon
        \right)
        \longrightarrow0.
    \]
    This proves \eqref{eq:mean-field-classification-error}.
\end{enumerate}
\end{proof}
